% Translated from the unchanged Typst manuscript; hardness proofs moved to Appendix D.

\section{Introduction}
Making sure that an artificial superintelligence does not pursue goals that are unaligned with humans is a challenging problem, and arguably one of the most important of our times~\citep{ngo2024alignment,bengio2024managing}. Several attempts have been made to get a handle on it and a recurring theme is that at its core it is a learning problem~\citep{hadfieldmenell2016cooperative,hadfieldmenell2017inverse}. Indeed, imagine several agents coexisting in the universe and interacting with each other with each agent trying to act in the interest of a subset of the agents. The agents only have partial information about each other's behaviours and goals as well as about the universe and must ensure the universe evolves towards desirable states~\citep{dafoe2020cooperative}. Given the role of learning, it is not unreasonable to imagine that learning theory might help us pave the way towards solving the problem~\citep{kosoy2018learningagenda}.

However, the currently known results in learning theory are likely insufficient to solve alignment. For example, it has been suggested that designing better learners for unrealizable learning is crucial for solving the problem~\citep{kosoy2018learningagenda}. Realizability is a strong assumption in general, but it becomes even more questionable in the multiagent learning scenarios described above.

One concrete issue with realizable learning is the grain of truth problem~\citep{hutter2009open,leike2016grain,wyeth2025grain,kosoy2023agenda}. Suppose Alice and Bob are playing a known game. Neither knows the other's strategy but assumes that they lie in a class $H$ of strategies~\citep{kalai1993rational,noguchi2015bayesian}. Can each player simultaneously achieve a low regret? The issue is that the strategy that achieves a low regret against all strategies in $H$ may not itself be a member of $H$~\citep{nachbar1997prediction,nachbar2005beliefs}. In that case, if Alice's strategy is in $H$, then Bob's strategy will have to be outside of $H$ and at least one of the players need to solve an unrealizable learning problem. Finding a general and natural class $H$ whose learner lies within $H$ is an open problem.






\section{Setting}
An instance of the imprecise bandits problem is a tuple $\mathcal{M} = \left( X , D , H , r \right)$, where $X$ is a set of arms, $D$ is a set of possible outcomes, $H$ is a set of mappings of type $X \to \credal D$ where $\credal D$ denotes the set of credal sets over $D$, and $r : X \times D \to \mathbb{R}$ is a reward function.

The imprecise bandits game is played between a learner and an adversary for $T$ rounds (both adversary and learner know $T$). At the beginning of the game, the adversary picks some $h^{\star} \in H$. The learner knows $X , D , H , r$, and $T$, but does not have knowledge of $h^{\star}$. Then, at each round $t$:

\begin{enumerate}
\item Learner picks an arm $x_{t} \in X$.

\item Adversary picks a distribution $P_{t} \in h^{\star} \left( x_{t} \right)$.

\item An outcome $y_{t} \sim P_{t}$ is drawn and shown to the learner as feedback.

\item Learner gets reward $r \left( x_{t} , y_{t} \right)$.

\end{enumerate}
Given a hypothesis $h \in H$, we define the value functions:


\begin{equation*}
v_{h} \left( x \right) \coloneqq \operatorname*{min}_{P \in h \left( x \right)} \mathbb{E}_{y \sim P} \left[ r \left( x , y \right) \right] , \quad V_{h} \coloneqq \operatorname*{max}_{x \in X} v_{h} \left( x \right) .
\end{equation*}


We use the shorthand $v^{\ast} \left( x \right)$ and $V^{\star}$ to denote the value functions corresponding to the true hypothesis $h^{\star}$.

The regret over horizon $T$ is


\begin{equation*}
R_{T} \coloneqq T V^{\star} - \mathbb{E} \left[ \sum_{t = 1}^{T} r \left( x_{t} , y_{t} \right) \right] .
\end{equation*}


Here the expectation in the second term is taken wrt the randomness in the learner's and the adversary's strategies as well as the randomness arising from sampling $y_{t} \sim P_{t}$.

We are interested in learner strategies that can achieve a sublinear (in $T$) regret for all true hypotheses $h^{\star} \in H$. We say that an imprecise bandits instance $\mathcal{M} = \left( X , D , H , r \right)$ has a \emph{statistically efficient} learner if such a learner strategy exists.

\textbf{Computational tractability.} In this paper, we are interested in studying the computational complexity of the learner. A \emph{learner policy} $\pi$ is a randomized algorithm that takes as input a history $h_{t} = {\left( \left( x_{i} , y_{i} \right) \right)}_{i = 1}^{t}$ and the description of the problem instance $\mathcal{M}$ and returns a next arm $x_{t + 1}$. The computational complexity of the problem depends crucially on the exact nature of access the policy is given to $\mathcal{M}$. For example, on one extreme, one can consider a non-uniform access such that a policy $\pi^{\mathcal{M}}$ is merely indexed by the instance $\mathcal{M}$ as opposed to $\mathcal{M}$ being an input to a computational procedure. However, this makes several instances rather trivial. In this paper we mostly consider specific encodings of $\mathcal{M}$. We also occasionally assume $\pi$ to have oracle access to $\mathcal{M}$ with certain specific oracles. We will make these assumptions clear in the respective sections.

We will say that a learner is \emph{computationally efficient} if it runs in time polynomial in the size of its input and achieves a regret $R_{T} \le \operatorname{poly} \left( | \mathcal{M} | \right) T^{1 - \alpha}$ for some $\alpha > 0$.


\subsection{Linear Imprecise Bandits}

For most of this work, we study the following linear specialization of the
instance $\mathcal{M} = \left( X , D , H , r \right)$ defined above. Let
$X \subseteq \mathbb{R}^{d_{X}}$ be the arm set, and let
$D \subseteq \mathbb{R}^{d_{D}}$ be the outcome set. We assume that $X$ and $D$ are
Euclidean balls of radii $R_{X}$ and $R_{D}$ (not necessarily centred at the origin).
The reward is affine in the arm and outcome:


\begin{equation*}
r \left( x , y \right) \coloneqq a^{T} x + b^{T} y + c .
\end{equation*}


Let the hypothesis class $H$ be parametrized by $Z \subseteq \mathbb{R}^{d_{Z}}$.
For every $z \in Z$, we define the hypothesis $h_{z}$ by first defining the possible expected values of the distributions that $h_{z}$ is allowed to choose. In particular, let



\begin{equation}
K_{z} \left( x \right) \coloneqq \left\{ y \in D : C_{z} y + B_{z} x + d_{z} = 0 \right\} .\label{eq:setting-compatible-set}
\end{equation}



\draft{
Here $B_{z}$, $C_{z}$, and $d_{z}$ are matrices and vectors of appropriate
dimensions that depend affinely on $z$. Now, hypothesis
$h_{z} : X \to \credal D$ is given by
}


\begin{equation*}
h_{z} \left( x \right) \coloneqq \left\{ P \in \Delta D : \mathbb{E}_{y \sim P} \left[ y \right] \in K_{z} \left( x \right) \right\} .
\end{equation*}


Let $z^{\star} \in Z$ represent the true
hypothesis, so that $h^{\star} = h_{z^{\star}}$. Define
$K^{\star} \left( x \right) \coloneqq K_{z^{\star}} \left( x \right)$.

We also assume a transversality condition as in~\citep{kosoy2025imprecise}: there is a constant
$S \in \left( 0 , 1 \right]$ such that, for every $x \in X$ and every
$p \in \mathbb{R}^{d_{D}}$ satisfying
$C_{z^{\star}} p + B_{z^{\star}} x + d_{z^{\star}} = 0$ with $p \notin D$,


\begin{equation*}
\operatorname{dist} \left( p , D \right) \ge S \operatorname{dist} \left( p , K^{\star} \left( x \right) \right) .
\end{equation*}


For a Euclidean ball, this says that the affine solution spaces defined by
the true constraints remain uniformly bounded away from tangency to $D$.

Because $r$ is affine in its outcome argument, the value functions defined
above satisfy


\begin{equation*}
v^{\ast} \left( x \right) = \operatorname*{min}_{P \in h^{\star} \left( x \right)} \mathbb{E}_{y \sim P} \left[ r \left( x , y \right) \right] = \operatorname*{min}_{y \in K^{\star} \left( x \right)} r \left( x , y \right) , \quad V^{\star} \coloneqq \operatorname*{max}_{x \in X} v^{\ast} \left( x \right) .
\end{equation*}


The regret over horizon $T$ is


\begin{equation*}
R_{T} \coloneqq T V^{\star} - \mathbb{E} \left[ \sum_{t = 1}^{T} r \left( x_{t} , y_{t} \right) \right] .
\end{equation*}


To define the computational problem, we assume that the numerical data defining $\mathcal{M}$ are rational and
given explicitly, with total encoding length
$|\mathcal{M}|\le\operatorname{poly}(d_X,d_D,d_Z)$.
Running time is measured in bit operations.

Now we state our main result.

\draft{
\begin{theorem}[efficient $\sqrt{T}$ learning]\label{thm:efficient-upper-bound}

  There exists a policy for linear imprecise bandits as defined above that,
  for every time horizon $T$ and every true hypothesis $h^{\star} \in H$,
  runs in time polynomial in $| \mathcal{M} |$ and $T$ and achieves regret

  
\begin{equation*}
R_{T} \le \tilde{O} \left( \sqrt{T} \right) .
\end{equation*}


  The hidden factor depends polynomially on the dimensions, $S^{- 1}$,
  and the reward range. The precise dependence is given in
  Theorem~\ref{thm:soft-square-root-regret}.

\end{theorem}


This dependence on $T$ is tight up to logarithmic factors: we can reuse
the $\Omega \left( \sqrt{T} \right)$ worst-case lower bound for stochastic bandits
\citep[Theorem 5.1]{auer2002nonstochastic}. To see this, take
$X = \left[ - 1 , 1 \right]$, $D = \left[ 0 , 1 \right]$, $r \left( x , y \right) = y$, and, for unknown
$\mu = \left( \mu_{1} , \mu_{2} \right) \in {\left[ 0 , 1 \right]}^{2}$, set


\begin{equation*}
K_{\mu} \left( x \right) = \frac{\{ \left( 1 - x \right)}{2} \mu_{1} + \frac{1 + x}{2} \mu_{2} \} .
\end{equation*}


These are affine constraints of the required form, with $S = 1$.
Given a two-armed Bernoulli bandit with means $\mu_{1}$ and $\mu_{2}$, we can
simulate a play of $x$ by pulling its first arm with probability
$\frac{1 - x}{2}$ and its second arm otherwise, and returning the observed reward
as $y$. This gives a compatible outcome distribution at $x$ and preserves
expected regret, since $V^{\star} = \operatorname*{max} \left( \mu_{1} , \mu_{2} \right)$. Thus a better worst-case
regret bound here would give a better bound for two-armed stochastic bandits.
}

{

\section{Warmup: A $T^{\frac{2}{3}}$ Learner}\label{sec:hard-ellipsoid-warmup}

For an arm $x$ and outcome $y$, consider the vector $w \left( x , y \right) = \left( 1 , x , y \right)$.
This vector lies in $\mathbb{R}^{1 + d_{X} + d_{D}}$, and the constraints defining the
feasible set of outcomes, $C_{z^{\star}} y + B_{z^{\star}} x + d_{z^{\star}} = 0$,
determine a linear subspace of this vector space. Define $\mathcal{L} \coloneqq \operatorname{ker} \begin{pmatrix}d_{z^{\star}} & B_{z^{\star}} & C_{z^{\star}}\end{pmatrix} .$ A conditional mean $m$ is feasible at $x$ precisely when
$m \in D$ and $w \left( x , m \right) \in \mathcal{L}$.
}{The learner's task is to learn enough about $\mathcal{L}$ to get low regret.
Of course, we may never learn all of $\mathcal{L}$. For example, if nature keeps
choosing conditional means in a strictly smaller subspace, then we cannot
distinguish that subspace from the true $\mathcal{L}$. But this is fine, since we do
not need to learn directions that nature never uses.}

The fact that we never get to observe the true conditional mean, but only a noisy observation of it, adds an extra complication. In a noiseless world, a new observation tells us a new direction of $\mathcal{L}$. Thus we can simply maintain a linear subspace of $\mathbb{R}^{1 + d_{Z} + d_{D}}$, and incrementally grow it every time we get new information. However, in the presence of noise, no observation can be used to add an entire direction with certainty. Instead, we encode our knowledge of $\mathcal{L}$ with an \emph{ellipsoid}, and incrementally grow it as we get new information. The details are written below.

{
We first translate and rescale the balls so that
$X = \left\{ x : {\left\Vert x \right\Vert}_{2} \le 1 \right\}$ and $D = \left\{ y : {\left\Vert y \right\Vert}_{2} \le 1 \right\}$. We continue to write the
reward as $r \left( x , y \right) = a^{T} x + b^{T} y + c$ and the true subspace as $\mathcal{L}$ in
these coordinates. The new ${\left\Vert b \right\Vert}_{2}$ equals $R_{D} {\left\Vert b \right\Vert}_{2}$ in the
original coordinates, and the non-tangency constant $S$ stays the same.
Write $q \coloneqq 1 + d_{X} + d_{D}$ for the number of coordinates of $w \left( x , y \right)$.
For a positive-definite matrix $V$, write
${\left\Vert w \right\Vert}_{V^{- 1}} \coloneqq \sqrt{w^{T} V^{- 1} w}$.


Define $p_{V} \left( x , y \right) \coloneqq {\left\Vert w \left( x , y \right) \right\Vert}_{V^{- 1}}^{2} ,$ and $E_{V} \coloneqq \left\{ w \in \mathbb{R}^{q} : {\left\Vert w \right\Vert}_{V^{- 1}}^{2} \le 1 \right\} .$


Here $E_{V}$ is the ellipsoid we described above, which we incrementally
expand as we get new information about $\mathcal{L}$.
The algorithm works in blocks of length $n$ (we will choose $n$ later
to get a small regret). At the beginning of each block, we choose an
arm $x$ and play the same arm throughout the block.
The ellipsoid $E_{V}$ stays fixed during these rounds.

To choose an arm, we use our current guess for the feasible conditional
means. For each $x$, this consists of the $m \in D$ with
$\left( 1 , x , m \right) \in E_{V}$. If some arm has no such $m$, we choose it to learn
more about it. Otherwise, we choose the arm with the largest worst-case
reward over its guessed feasible means.

At the end of a full block, we average its $n$ observations to get
$\overline{y}$. We then check whether $\left( 1 , x , \overline{y} \right)$ lies in $E_{V}$.
If it does not, we enlarge the ellipsoid to include it. Otherwise, we
leave the ellipsoid unchanged.

For the time being, we assume that we can carry out the rule for
choosing an arm exactly. We will consider the computational aspects
of the algorithm later.


\begin{algorithm}
\floatconts{alg:hard-ellipsoid-warmup}{\caption{}}{

  Fix a block length $n \in \left\{ 1 , \ldots , T \right\}$ and initialize $V \coloneqq n^{- 1} I_{q}$.
  At the start of each block:


\begin{enumerate}
\item For each arm, let

    
\begin{equation*}
{\hat{K}}_{V} \left( x \right) \coloneqq \left\{ y \in D : p_{V} \left( x , y \right) \le 1 \right\} .
\end{equation*}


    If any of these sets is empty, choose an arm $x$ with
    ${\hat{K}}_{V} \left( x \right) = \varnothing$. Otherwise, define

    
\begin{equation*}
H_{V} \left( x \right) \coloneqq \operatorname*{min}_{y \in {\hat{K}}_{V} \left( x \right)} r \left( x , y \right) ,
\end{equation*}


    and choose

    
\begin{equation*}
x \in \operatorname{argmax}_{x^{\prime} \in X} H_{V} \left( x^{\prime} \right) .
\end{equation*}

\item Play $x$ for $n$ rounds, or until the horizon if fewer rounds remain.
    Let $\overline{y}$ be the empirical mean of observed outcomes and
    $\overline{w} \coloneqq w \left( x , \overline{y} \right)$.

\item If the block has length $n$ and $p_{V} \left( x , \overline{y} \right) > 1$, update
    $V \leftarrow V + \overline{w} {\overline{w}}^{T}$. Otherwise leave $V$ unchanged.

\end{enumerate}

}
\end{algorithm}
\FloatBarrier



\draft{
\begin{theorem}[Informal]\label{thm:hard-ellipsoid-warmup}

  For fixed problem parameters, Algorithm~\ref{alg:hard-ellipsoid-warmup} with exact
  planning has expected regret

  
\begin{equation*}
R_{T} \le \tilde{O} \left( n + \frac{T}{\sqrt{n}} \right) .
\end{equation*}


  In particular, choosing $n = \left\lceil T^{\frac{2}{3}} \right\rceil$ gives
  $R_{T} \le \tilde{O} \left( T^{\frac{2}{3}} \right)$.

\end{theorem}
}

}

{
\begingroup
\renewcommand{\proofname}{Proof sketch}
\begin{proof}
We will show that, with high probability, throughout the
$T$ time steps, every point in $E_{V}$ is at most about $\frac{1}{\sqrt{n}}$ away
from the true $\mathcal{L}$ (Lemma~\ref{lem:subspace-distance-bound}). We ignore
constants depending on the fixed problem parameters and logarithmic
factors in this sketch.

Now consider an uninformative block, meaning a full block in which we
do not update $V$. The arm $x$ that we play during this block is picked
by maximizing $H_{V}$. Since the block is uninformative, its
empirical mean outcome satisfies $w \left( x , \overline{y} \right) \in E_{V}$. These two
facts, together with the distance bound and Lemma~\ref{lem:subspace-reward-bound},
imply that the average reward is at least $V^{\star}$ minus about
$\frac{1}{\sqrt{n}}$. Thus the block's regret (the best arm's guaranteed reward
$n V^{\star}$ minus the reward we actually achieved) is at most about
$\frac{n}{\sqrt{n}} = \sqrt{n}$.

On the other hand, for an informative block, the empirical mean can
lie outside the ellipsoid, so we only have the trivial regret bound
of $O \left( n \right)$.

Finally, the update rule for $V$ ensures that there are at most
$O \left( \operatorname{log} T \right)$ informative blocks: each update more than doubles the
determinant, whose total growth is bounded by a polynomial in $T$
(see the proof of Theorem~\ref{thm:hard-ellipsoid-regret}). There are at most $\frac{T}{n}$
uninformative blocks, so the total regret is about $n \operatorname{log} T + \sqrt{n} \frac{T}{n}$.

Ignoring logarithmic factors, this bound is minimized by choosing
$n = \left\lceil T^{\frac{2}{3}} \right\rceil$, giving an upper bound of
$\tilde{O} \left( T^{\frac{2}{3}} \right)$ on the regret.
\end{proof}
\endgroup

The precise statement and proof are in Appendix~\ref{app:hard-ellipsoid-analysis}.
}

\subsection{Computational tractability}\label{sec:hard-computation}
In this section we describe a polytime implementation of Algorithm~\ref{alg:hard-ellipsoid-warmup}.

The matrix updates $V \leftarrow V + \overline{w}\overline{w}^{T}$ and the test $p_{V} \left( x , \overline{y} \right) > 1$ can be computed
in polynomial time. The main issue is choosing an arm, which involves solving the following two computational problems:

\begin{enumerate}
\item Decide whether there exists an arm $x\in X$ with
  $\hat{K}_V(x)=\varnothing$, and if so, find such an arm.
\item If $\hat{K}_V(x)$ is nonempty for every $x\in X$, find an arm
  $x\in\operatorname*{argmax}_{x'\in X}H_V(x')$, where
  $H_V(x)=\min_{y\in\hat{K}_V(x)}r(x,y)$.
\end{enumerate}

We use the following result of Bienstock~\citep[Theorem~1.3]{bienstock2016cdt} for both problems.

\begin{theorem}[Bienstock~\citep{bienstock2016cdt}]
  \label{thm:bienstock}
Consider a maximization problem with a nonempty feasible set that has: a) a fixed number of quadratic constraints with at least one of them being a bounded ellipsoid, b) a quadratic objective function, and c) rational coefficients for both the objective function and the constraints. Then, there exists an algorithm that, for every
$\varepsilon\in(0,1)$, returns rational assignments to the variables that violate
each constraint by at most $\varepsilon$ and achieve an objective value that is at
least the optimum minus $\varepsilon$. Its running time is polynomial in
the input size and $\log(1/\varepsilon)$.
\end{theorem}

Since all our constraints and objective functions are quadratic, one can show that both problems above can be written in a way that Theorem~\ref{thm:bienstock} can be applied. One can choose the $\varepsilon$ as an appropriate function of $T$ such that the approximate solution is still good enough that it does not increase the regret. Details are in Appendix~\ref{app:hard-computation}.

\section{A $\sqrt{T}$ learner}\label{sec:soft-square-root}

{
The previous section attempts to apply the principle of ``optimism in the
face of uncertainty,'' but does not fully embrace it. Indeed, the arm for
each block is picked optimistically, but each block still looks like a
standard exploration block. The learner does not always need to wait for
$n$ steps before identifying that the arm it is playing is problematic.
If the outcomes collected so far have an empirical mean $\overline{y}$
with $w(x,\overline{y})$ already far outside the current ellipsoid, then
the block can be cut short more quickly. However, the current arm score $H_V(x)$ 
does not account for outcomes outside the ellipsoid. First, we design a
score function that takes every outcome in $D$ into account, not just
those represented inside our chosen ellipsoid. The score function will
be linked to the regret in a way that lets us decide when to move on to
the next block.
}

{
We construct this score by replacing the hard constraint
$p_{V} \left( x , y \right) \le 1$ in $H_{V}$ with a penalty, so that we can
minimize over all of $D$. For a parameter $\lambda > 0$, which we will
choose later to get a small regret, define


\begin{equation*}
U_{V} \left( x \right) \coloneqq \operatorname*{min}_{y \in D} \left[ r \left( x , y \right) + \lambda p_{V} \left( x , y \right) \right] .
\end{equation*}


We pick the arm $x$ by maximizing $U_V$. For now, assume that we can
maximize the score exactly.

\noindent The following lemma gives us a way to control regret when the outcome $y$
appears anywhere in $D$.

\begin{lemma}\label{lem:soft-outcome-regret}
  Fix $\delta\in(0,1)$. With probability at least $1-\delta$, simultaneously
  for every matrix $V$ that Algorithm~\ref{alg:soft-square-root} considers,
  every $\lambda>0$, and every arm $x$ maximizing $U_V$ with this value of $\lambda$,

\begin{equation*}
V^\star-r(x,y)
\le \lambda p_V(x,y)+\tilde{O}(1/\lambda)
\qquad\text{for every }y\in D.
\end{equation*}

  Here $\tilde{O}$ hides logarithmic factors in $T/\delta$.

\end{lemma}

\draft{
\begin{proof}
On the event of Lemma~\ref{lem:interval-noise}, the distance
bound in Equation~\ref{eq:soft-subspace-certificate}, proved in
Appendix~\ref{app:soft-square-root-analysis}, and
Lemma~\ref{lem:subspace-reward-bound} give, for every arm $x'$ and
outcome $y\in D$,

\begin{equation*}
r(x',y)\ge v^\star(x')-\tilde{O}\left(\sqrt{p_V(x',y)}\right).
\end{equation*}

The penalty grows quadratically in $\sqrt{p_V(x',y)}$, while the error
bound grows linearly. In particular,
$\lambda p-\tilde{O}(\sqrt p)\ge-\tilde{O}(1/\lambda)$ for every $p\ge0$
and every $\lambda>0$.
Adding the penalty and minimizing over $y$ therefore gives
$U_V(x')\ge v^\star(x')-\tilde{O}(1/\lambda)$.
Since $x$ maximizes the score,
$U_V(x)\ge V^\star-\tilde{O}(1/\lambda)$.
Finally, the definition of the score gives
$U_V(x)\le r(x,y)+\lambda p_V(x,y)$ for every $y\in D$.
Combining these bounds proves the claim. Since the distance bound is
independent of $\lambda$, the argument holds for all $\lambda>0$ on the
same event.
\end{proof}
}

To use this bound during a block, we keep $V$ fixed and play the same
arm $x$ repeatedly, but recompute the empirical mean after every
observation. Let ${\overline{y}}_n$ be the mean of the first $n$ outcomes
in the block. Since the reward is affine, the reward accumulated so far
is $n r(x,{\overline{y}}_n)$. Applying Lemma~\ref{lem:soft-outcome-regret}
to this mean gives

\begin{equation*}
n V^\star-\sum_{i=1}^n r(x,y_i)
\le \lambda n p_V(x,{\overline{y}}_n)+\tilde{O}(n/\lambda).
\end{equation*}

The first term is observable. We can compute it after every observation,
and we end the block and update $V$ the first time it reaches $\lambda$,
or equivalently, when $n p_V(x,{\overline{y}}_n)\ge1$.
Before this happens, the block's regret is bounded by
$\lambda+\tilde{O}(n/\lambda)$.

A large empirical penalty can therefore make us update after only a
few observations. A small one lets us keep playing the arm for longer. Thus we choose when to update $V$ by monitoring this bound
throughout the block, rather than waiting for a fixed number of rounds.

\begin{algorithm}[!htbp]
{
\floatconts{alg:soft-square-root}{\caption{}}{

  Initialize $V \coloneqq I_{q}$. At the start of each block:


\begin{enumerate}
\item Define

    
\begin{equation*}
U_{V} \left( x \right) \coloneqq \operatorname*{min}_{y \in D} \left[ r \left( x , y \right) + \lambda p_{V} \left( x , y \right) \right] ,
\end{equation*}


    and choose an arm
    $x \in \operatorname*{argmax}_{x^{\prime} \in X} U_{V} \left( x^{\prime} \right)$.

\item Play $x$ repeatedly. After $n$ observations in the block, let

    
\begin{equation*}
{\overline{y}}_{n} \coloneqq \frac{1}{n} \sum_{i = 1}^{n} y_{i} , \quad {\overline{w}}_{n} \coloneqq w \left( x , {\overline{y}}_{n} \right) .
\end{equation*}


    End the block the first time

    
\begin{equation*}
n p_{V} \left( x , {\overline{y}}_{n} \right) \ge 1 .
\end{equation*}

\item Update $V \leftarrow V + n {\overline{w}}_{n} {\overline{w}}_{n}^{T}$ and start a new block.
    At the horizon, stop even if the current block has not ended.

\end{enumerate}

}
}
\end{algorithm}
% \FloatBarrier

\begin{theorem}[Informal]\label{thm:soft-square-root}

  For fixed problem parameters,
  Algorithm~\ref{alg:soft-square-root}, with $\lambda = \lceil \sqrt{T} \rceil$,
  has expected regret

\begin{equation*}
R_T \le \tilde{O}\left(\sqrt{T}\right).
\end{equation*}

  The learner does not need to know $S$.

\end{theorem}

\begingroup
\renewcommand{\proofname}{Proof sketch}
\draft{
\begin{proof}
As before, each completed block at least doubles
the determinant of $V$, so there are
only $O(\log T)$ blocks.

By Lemma~\ref{lem:soft-outcome-regret}, with high probability, a block of
length $n$ has regret at most
$\lambda n p_V(x,\overline{y})+\tilde{O}(n/\lambda)$.
The stopping rule keeps $n p_V(x,\overline{y})$ bounded by a constant,
so the first term costs $O(\lambda)$ per block, regardless of its length.
The remaining term costs about $1/\lambda$ per round.
Summing both costs gives regret
at most about $\lambda \log T + T/\lambda$. Ignoring
logarithmic factors, choosing $\lambda = \lceil\sqrt{T}\rceil$ balances
the two terms. Taking the concentration failure probability of
order $T^{-2}$ gives the same $\tilde{O}(\sqrt{T})$ bound in expectation.
\end{proof}
}
\endgroup

The precise regret bound and detailed proof are in
Appendix~\ref{app:soft-square-root-analysis}.

\draft{
\subsection{Computational tractability}\label{sec:soft-computation}

To implement the learner, we allow an additive error $\varepsilon>0$
when maximizing $U_V$. This adds at most $n\varepsilon$ to the regret
of a block of length $n$, and at most $T\varepsilon$ over all rounds.
Taking $\varepsilon=(T+1)^{-2}$ therefore preserves the
$\tilde{O}(\sqrt{T})$ regret bound.

Dualizing the penalty turns the planning step into a quadratic program
with a fixed number of quadratic constraints. The weak-optimization
result in Appendix~\ref{app:qcqp} and sufficiently fine rounding give
an implementation under the weak finite-precision model that runs in
time polynomial in $|\mathcal{M}|$ and $T$ and retains the same regret
rate. The details are in Appendix~\ref{app:soft-square-root-analysis}.

}


\section{NP-hardness}

\draft{Proofs are in Appendix~\ref{app:np-hardness-proofs}.}

\subsection{D = simplex, X = ball}

\begin{theorem}

  Consider a variation of the linear imprecise bandits problem where $D$ is allowed to be a simplex. If, for
  some fixed $\beta > 0$, a randomized polynomial-time learner guarantees $E \left[ R_{T} \right] \le P \left( | \mathcal{M} | \right) T^{1 - \beta}$ for every instance $\mathcal{M}$, then $\text{\normalfont NP} \subseteq \text{\normalfont BPP}$. If a deterministic learner achieves the same regret guarantee, then $\text{\normalfont P} = \text{\normalfont NP}$.

 \label{thm:additive-simplex-np-hardness}

\end{theorem}

\draft{

\subsection{A Euclidean outcome ball and a polytope of arms}
}

\begin{theorem}\label{thm:additive-polytope-np-hardness}

  \draft{
  Suppose we allow $X$ to be a polytope in the linear imprecise bandits
  problem. If, for some fixed $\beta > 0$ and polynomial $P$, a randomized
  polynomial-time learner guarantees $R_{T} \le P \left( | \mathcal{M} | \right) T^{1 - \beta}$ for every
  instance $\mathcal{M}$, then
  $\text{\normalfont NP} \subseteq \text{\normalfont BPP}$. If a deterministic learner achieves the same
  regret guarantee, then $\text{\normalfont P} = \text{\normalfont NP}$.
  }

\end{theorem}

\draft{

\subsection{Trilinear constraints with Euclidean arm and outcome balls}
}

\begin{theorem}\label{thm:trilinear-balls-np-hardness}
  \draft{
  For the instances below,
  $F:\mathbb{R}^{d_X}\times\mathbb{R}^{d_Z}\times Y\to W$
  is linear in each argument separately, where $Y$ is the ambient outcome
  space and $W$ is the constraint space. Write
  $F_{x,z}(y):=F(x,z,y)$, $L_z(x):=\ker F_{x,z}$, and
  $K_z(x):=L_z(x)\cap D$, where $D$ lies in the hyperplane
  in which an added outcome coordinate equals one.
  }


  \draft{
  There is a polynomial-time reduction from CLIQUE to imprecise-bandit
  instances with rational coefficients for which $X$ is a full-dimensional
  Euclidean ball, $D$ is a Euclidean ball in its affine hull, $Z = \left[ 1 , 2 \right]$ is a
  one-dimensional Euclidean ball, the reward is linear, and $F$ is
  trilinear. Every set $K_{z} \left( x \right)$ is a singleton independent of $z$.
  The non-tangency inequality holds uniformly for every point in
  $L_{z} \left( x \right)$ outside $D$, with a constant $S > 0.88$.
  }

  \draft{
  For these instances, approximating $V_{z}$ to inverse-polynomial additive
  accuracy is NP-hard even when $z$ is known. Thus, if, for some
  fixed $\beta > 0$, a randomized polynomial-time learner guaranteed
  }

  \draft{

\begin{equation*}
R_{T} \le P \left( | \mathcal{M} | \right) T^{1 - \beta}
\end{equation*}

  }

  \draft{
  on every such instance for a polynomial $P$, then $\text{\normalfont NP} \subseteq \text{\normalfont BPP}$.
  For a deterministic learner, the same conclusion would be $\text{\normalfont P} = \text{\normalfont NP}$.
  }

\end{theorem}

\draft{

\subsection{Computing IUCB's first arm is NP-hard even for Euclidean balls}
}

\begin{theorem}\label{thm:iucb-np-hardness}
  \draft{
  Write $V_{\mathrm{IUCB}}:=\max_{z\in Z,\,x\in X}v_z(x)$ for
  IUCB's optimistic value on its first round. A globally optimal
  first arm maximizes $\max_{z\in Z}v_z(x)$ over $x\in X$.
  }


  \draft{
  Computing $V_{\text{\normalfont IUCB}}$ exactly or finding a globally optimal first arm
  of IUCB is NP-hard, even when $X$, $D$, and $Z$ are Euclidean balls and
  planning for each fixed hypothesis takes polynomial time.
  }

\end{theorem}

\draft{

\subsection{Efficient planning does not imply efficient learning}
}

\begin{theorem}\label{thm:easy-planning-hard-learning}

  \draft{
  There is a family of rational imprecise-bandit instances indexed by
  integers $2 \le k \le n$, with descriptions of length polynomial in $n$.
  Each instance has finite arm and hypothesis sets, an outcome set that is a
  rational polytope, constraints $F = F_{0} + F_{1}$ with $F_{0} = 0$ and $F_{1}$ bilinear
  in the outcome and hypothesis, and a reward
  in $\left[ 0 , 1 \right]$ that is linear in the arm and in the outcome separately.
  An exact planner takes $O \left( n^{2} \right)$ time when the hypothesis is known.
  }

  \draft{
  Suppose that, for some fixed $\beta > 0$ and polynomial $P$, one uniform
  randomized policy runs in time polynomial in the description length
  $| \mathcal{M} |$ and horizon $N$ and has expected regret at most
  }

  \draft{

\begin{equation*}
P \left( | \mathcal{M} | \right) N^{1 - \beta}
\end{equation*}

  }

  \draft{
  for every instance in this family, every true hypothesis, and every
  compatible nature policy. Then $\text{\normalfont CLIQUE} \in \text{\normalfont RP}$ and hence
  $\text{\normalfont NP} = \text{\normalfont RP}$. The same conclusion holds if the guarantee is required
  only against stationary deterministic nature policies. A deterministic
  learner with the same guarantee would imply $\text{\normalfont P} = \text{\normalfont NP}$.
  }

\end{theorem}

% Rendered wording and order transcribed from the current Typst PDF.
% Keep IEEE references here to preserve the source document's visible content.
\begingroup
\color{black}
\renewcommand{\refname}{\draft{References}}
\begin{thebibliography}{25}
\bibitem{kosoy2025imprecise}
V. Kosoy, ``Imprecise Multi-Armed Bandits: Representing Irreducible Uncertainty as a Zero-Sum Game,'' \emph{Journal of Machine Learning Research}, vol. 26, no. 184, pp. 1--75, 2025, [Online]. Available: \url{https://www.jmlr.org/papers/v26/24-2001.html}

\bibitem{auer2002nonstochastic}
P. Auer, N. Cesa-Bianchi, Y. Freund, and R. E. Schapire, ``The Nonstochastic Multiarmed Bandit Problem,'' \emph{SIAM Journal on Computing}, vol. 32, no. 1, pp. 48--77, 2002, doi: \nolinkurl{10.1137/S0097539701398375}.

\bibitem{bienstock2016cdt}
D. Bienstock, ``A Note on Polynomial Solvability of the CDT Problem,'' \emph{SIAM Journal on Optimization}, vol. 26, no. 1, pp. 488--498, 2016, doi: \nolinkurl{10.1137/15M1009871}.

\bibitem{motzkin1965maxima}
T. S. Motzkin and E. G. Straus, ``Maxima for Graphs and a New Proof of a Theorem of Tur\'{a}n,'' \emph{Canadian Journal of Mathematics}, vol. 17, pp. 533--540, 1965, doi: \nolinkurl{10.4153/CJM-1965-053-6}.

\bibitem{hillar2013most}
C. J. Hillar and L.-H. Lim, ``Most Tensor Problems Are NP-Hard,'' \emph{Journal of the ACM}, vol. 60, no. 6, p. 45:1--45:39, 2013, doi: \nolinkurl{10.1145/2512329}.

\bibitem{garey1979computers}
M. R. Garey and D. S. Johnson, \emph{Computers and Intractability: A Guide to the Theory of NP-Completeness}. W. H. Freeman, 1979.

\bibitem{boyd2004convex}
S. Boyd and L. Vandenberghe, \emph{Convex Optimization}. Cambridge University Press, 2004. doi: \nolinkurl{10.1017/CBO9780511804441}.

\bibitem{pardalos1991negative}
P. M. Pardalos and S. A. Vavasis, ``Quadratic Programming with One Negative Eigenvalue Is NP-Hard,'' \emph{Journal of Global Optimization}, vol. 1, no. 1, pp. 15--22, 1991, doi: \nolinkurl{10.1007/BF00120662}.

\bibitem{barvinok1993feasibility}
A. I. Barvinok, ``Feasibility Testing for Systems of Real Quadratic Equations,'' \emph{Discrete \& Computational Geometry}, vol. 10, pp. 1--13, 1993, doi: \nolinkurl{10.1007/BF02573959}.

\bibitem{karp1972reducibility}
R. M. Karp, ``Reducibility Among Combinatorial Problems,'' in \emph{Complexity of Computer Computations}, R. E. Miller, J. W. Thatcher, and J. D. Bohlinger, Eds., New York: Plenum Press, 1972, pp. 85--103. doi: \nolinkurl{10.1007/978-1-4684-2001-2_9}.

\bibitem{garey1976simplified}
M. R. Garey, D. S. Johnson, and L. Stockmeyer, ``Some Simplified NP-Complete Graph Problems,'' \emph{Theoretical Computer Science}, vol. 1, no. 3, pp. 237--267, 1976, doi: \nolinkurl{10.1016/0304-3975(76)90059-1}.

\bibitem{ngo2024alignment}
R. Ngo, L. Chan, and S. Mindermann, ``The Alignment Problem from a Deep Learning Perspective,'' in \emph{International Conference on Learning Representations}, 2024, [Online]. Available: \url{https://arxiv.org/abs/2209.00626v6}

\bibitem{bengio2024managing}
Y. Bengio \emph{et al.}, ``Managing extreme AI risks amid rapid progress,'' \emph{Science}, vol. 384, no. 6698, pp. 842--845, 2024, doi: \nolinkurl{10.1126/science.adn0117}.

\bibitem{hadfieldmenell2016cooperative}
D. Hadfield-Menell, A. Dragan, P. Abbeel, and S. Russell, ``Cooperative Inverse Reinforcement Learning,'' in \emph{Advances in Neural Information Processing Systems}, vol. 29, pp. 3909--3917, 2016, [Online]. Available: \url{https://arxiv.org/abs/1606.03137}

\bibitem{hadfieldmenell2017inverse}
D. Hadfield-Menell, S. Milli, P. Abbeel, S. Russell, and A. Dragan, ``Inverse Reward Design,'' in \emph{Advances in Neural Information Processing Systems}, vol. 30, 2017, [Online]. Available: \url{https://arxiv.org/abs/1711.02827}

\bibitem{dafoe2020cooperative}
A. Dafoe, E. Hughes, Y. Bachrach, T. Collins, K. R. McKee, J. Z. Leibo, K. Larson, and T. Graepel, ``Open Problems in Cooperative AI,'' arXiv preprint arXiv:2012.08630, 2020, [Online]. Available: \url{https://arxiv.org/abs/2012.08630}

\bibitem{kosoy2018learningagenda}
{\raggedright
V. Kosoy, ``The Learning-Theoretic AI Alignment Research Agenda,'' \emph{AI Alignment Forum}, Jul. 4, 2018, [Online]. Available: \url{https://www.alignmentforum.org/posts/5bd75cc58225bf0670375575/the-learning-theoretic-ai-alignment-research-agenda}\par}

\bibitem{hutter2009open}
M. Hutter, ``Open Problems in Universal Induction \& Intelligence,'' \emph{Algorithms}, vol. 2, no. 3, pp. 879--906, 2009, doi: \nolinkurl{10.3390/a2030879}.

\bibitem{leike2016grain}
J. Leike, J. Taylor, and B. Fallenstein, ``A Formal Solution to the Grain of Truth Problem,'' in \emph{Proceedings of the Thirty-Second Conference on Uncertainty in Artificial Intelligence}, pp. 427--436, 2016, [Online]. Available: \url{https://arxiv.org/abs/1609.05058}

\bibitem{wyeth2025grain}
C. Wyeth, M. Hutter, J. Leike, and J. Taylor, ``Limit-Computable Grains of Truth for Arbitrary Computable Extensive-Form (Un)Known Games,'' arXiv preprint arXiv:2508.16245, 2025, [Online]. Available: \url{https://arxiv.org/abs/2508.16245}

\bibitem{kosoy2023agenda}
{\raggedright
V. Kosoy, ``The Learning-Theoretic Agenda: Status 2023,'' \emph{AI Alignment Forum}, Apr. 19, 2023, [Online]. Available: \url{https://www.alignmentforum.org/posts/ZwshvqiqCvXPsZEct/the-learning-theoretic-agenda-status-2023}\par}

\bibitem{kalai1993rational}
E. Kalai and E. Lehrer, ``Rational Learning Leads to Nash Equilibrium,'' \emph{Econometrica}, vol. 61, no. 5, pp. 1019--1045, 1993, doi: \nolinkurl{10.2307/2951492}.

\bibitem{noguchi2015bayesian}
Y. Noguchi, ``Bayesian Learning, Smooth Approximate Optimal Behavior, and Convergence to $\varepsilon$-Nash Equilibrium,'' \emph{Econometrica}, vol. 83, no. 1, pp. 353--373, 2015, doi: \nolinkurl{10.3982/ECTA9332}.

\bibitem{nachbar1997prediction}
J. H. Nachbar, ``Prediction, Optimization, and Learning in Repeated Games,'' \emph{Econometrica}, vol. 65, no. 2, pp. 275--309, 1997, doi: \nolinkurl{10.2307/2171894}.

\bibitem{nachbar2005beliefs}
J. H. Nachbar, ``Beliefs in Repeated Games,'' \emph{Econometrica}, vol. 73, no. 2, pp. 459--480, 2005, doi: \nolinkurl{10.1111/j.1468-0262.2005.00585.x}.
\end{thebibliography}
\endgroup




\clearpage


\section*{\draft{Appendix}}


\appendix
\renewcommand{\presectionnum}{}
\renewcommand{\thesubsection}{\thesection.\Alph{subsection}}



\draft{

\section{Quadratically constrained quadratic programming}\label{app:qcqp}
}

\draft{
For symmetric matrices $Q_{i} \in \mathbb{R}^{d \times d}$, a \textbf{quadratically constrained
quadratic program} (QCQP) in $z \in \mathbb{R}^{d}$ has the form
}

\draft{

\begin{equation*}
\operatorname*{min}_{z \in \mathbb{R}^{d}} g_{0} \left( z \right) \quad \text{\normalfont subject to} \quad g_{i} \left( z \right) \le 0 , \quad i = 1 , \ldots , m ,
\end{equation*}

}

\draft{
where $g_{i} \left( z \right) \coloneqq z^{T} Q_{i} z + 2 p_{i}^{T} z + r_{i}$; for complexity statements, all
entries are rational. A quadratic equality $h \left( z \right) = 0$ is exactly the pair
$h \left( z \right) \le 0$ and $- h \left( z \right) \le 0$ \citep{bienstock2016cdt}.
}

\draft{
QCQP describes algebraic form, whereas convex programming describes geometry.
The displayed minimization problem is a convex program when every $Q_{i}$,
$i = 0 , \ldots , m$, is positive semidefinite and all equalities are affine
\citep{boyd2004convex}. Without those restrictions it can be nonconvex: general QCQP
is NP-hard even when only the objective is quadratic and its Hessian has one
negative eigenvalue \citep{pardalos1991negative}.
}

\draft{
The polynomial-time conclusion for our nonconvex instance has two steps.
}


\begin{enumerate}
\item \draft{
  First, fix the number $m$ of constraints. Bienstock shows that,
  for rational quadratics $g_{0} , \ldots , g_{m}$, if at least one constraint
  $g_{i} \left( z \right) \le 0$ has a positive-definite quadratic part, then for every
  $\varepsilon \in \left( 0 , 1 \right)$ an algorithm runs in time
  $\operatorname{poly} \left( L , \operatorname{log} \left( \frac{1}{\varepsilon} \right) \right)$, where $L$ is the input bit length, and either
  proves infeasibility or returns $\hat{z}$ such that
  }

  \draft{

\begin{equation*}
g_{i} \left( \hat{z} \right) \le \varepsilon , \quad i = 1 , \ldots , m ,
\end{equation*}

  }

  \draft{
  and $g_{0} \left( \hat{z} \right) \le g_{0} \left( z \right) + \varepsilon$ for every feasible $z$
  \citep[Theorem~1.3]{bienstock2016cdt}. The underlying weak-feasibility step reduces to a fixed
  number of homogeneous quadratic equations on a sphere, handled by
  Barvinok's method; binary search then gives the objective guarantee
  \citep{bienstock2016cdt}, \citep{barvinok1993feasibility}.
  }

\item \draft{
  Second, the planning reduction in Appendix~\ref{app:soft-square-root-analysis} has six
  quadratic inequalities. The redundant ellipsoid constraint on
  $\left( x , z , s \right)$ has a positive-definite quadratic part, so the QCQP
  satisfies the hypotheses of this result. Applying it to the negative
  objective computes the planning maximum to weak additive accuracy
  $\varepsilon$ in time $\operatorname{poly} \left( L , \operatorname{log} \left( \frac{1}{\varepsilon} \right) \right)$. The same appendix explains
  how to repair the solver's output to obtain a feasible arm. If either
  radius is zero, the problem first reduces to a lower-dimensional instance.
  }

\end{enumerate}
\draft{

\section{MAX-CUT}\label{app:max-cut}
}

\draft{
Let $G = \left( V , E \right)$ be a finite undirected simple graph. For $S \subseteq V$, the
edge boundary of $S$ is
}

\draft{

\begin{equation*}
\delta_{G} \left( S \right) \coloneqq \left\{ \left\{ u , v \right\} \in E : \left| \left\{ u , v \right\} \cap S \right| = 1 \right\} .
\end{equation*}

}

\draft{
The decision problem MAX-CUT takes as input $G$ and an integer
$k \in \left\{ 0 , \ldots , \left| E \right| \right\}$, and asks whether
}

\draft{

\begin{equation*}
\exists S \subseteq V \quad \text{\normalfont such that} \quad \left| \delta_{G} \left( S \right) \right| \ge k .
\end{equation*}

}

\draft{
A set $S$ satisfying this inequality is a polynomial-size certificate, so
MAX-CUT belongs to NP. Karp proved that weighted MAX-CUT is NP-complete
\citep{karp1972reducibility}. Garey, Johnson, and Stockmeyer subsequently proved
that the unweighted problem defined above remains NP-complete even when $G$
is cubic \citep{garey1976simplified}.
}

\draft{
The associated optimization problem computes
}

\draft{

\begin{equation*}
M_{G} \coloneqq \operatorname*{max}_{S \subseteq V} \left| \delta_{G} \left( S \right) \right| .
\end{equation*}

}

\draft{
Since the decision problem asks whether $M_{G} \ge k$, computing $M_{G}$ is NP-hard.
}

\draft{

\section{Regret analysis of the $T^{\frac{2}{3}}$ learner}\label{app:hard-ellipsoid-analysis}

We give the precise statement and proof of Theorem~\ref{thm:hard-ellipsoid-warmup}.
We use the coordinates and notation from Section~\ref{sec:hard-ellipsoid-warmup},
and assume that the planning steps in Algorithm~\ref{alg:hard-ellipsoid-warmup} are
solved exactly. A full block is \emph{informative} if it causes an update,
and \emph{uninformative} otherwise.


For the analysis, fix $\delta \in \left( 0 , 1 \right)$ and define


\begin{equation*}
\begin{aligned}
J_{\text{\normalfont max}} & \coloneqq \left\lceil q \operatorname{log}_{2} \left( 1 + 3 \frac{T}{q} \right) \right\rceil , \\
 h_{T}^{2} & \coloneqq 8 d_{D} \operatorname{log} \left( \frac{2 d_{D} T^{2}}{\delta} \right) , \\
 B_{T}^{2} & \coloneqq 1 + J_{\text{\normalfont max}} h_{T}^{2} , \\
 G & \coloneqq \sqrt{3} \left( 1 + S^{- 1} \right) {\left\Vert b \right\Vert}_{2} , \\
 C_{r} & \coloneqq \operatorname*{max}_{X \times D} r - \operatorname*{min}_{X \times D} r .
\end{aligned}
\end{equation*}



We write $J$ for the number of updates actually made. We will show
that $J \le J_{\text{\normalfont max}}$.

If $w \left( x , y \right)$ is close to $\mathcal{L}$, then $r \left( x , y \right)$ cannot be much smaller
than the worst feasible reward $v^{\star} \left( x \right)$. We first prove this.


\begin{lemma}\label{lem:subspace-reward-bound}

  For every $x \in X$ and $y \in D$,


\begin{equation*}
v^{\star} \left( x \right) \le r \left( x , y \right) + G \operatorname{dist} \left( w \left( x , y \right) , \mathcal{L} \right) .
\end{equation*}


\end{lemma}


\begin{proof}
Let $N \coloneqq \operatorname{ker} C_{z^{\star}}$ and let $Q$ be the orthogonal projector
onto $N^{\perp}$. Every arm has a feasible outcome, so we can write the
minimum-norm solution of the affine constraints as
$f \left( x \right) = A x + e \in N^{\perp}$. Thus


\begin{equation*}
K^{\star} \left( x \right) = \left( A x + e + N \right) \cap D .
\end{equation*}


Since there is a feasible outcome in the unit ball for every arm,
${\left\Vert A x + e \right\Vert}_{2} \le 1$ on $X$. Taking $x = 0$ gives ${\left\Vert e \right\Vert}_{2} \le 1$, and
comparing opposite unit vectors gives ${\left\Vert A \right\Vert}_{\text{\normalfont op}} \le 1$. Now define


\begin{equation*}
\Phi \left( s , x , y \right) \coloneqq Q y - A x - s e .
\end{equation*}


This map has operator norm at most $\sqrt{3}$ and kernel $\mathcal{L}$: the
original constraints are equivalent to $Q y = A x + e$, and homogenizing
either equation gives the same subspace. Therefore


\begin{equation*}
\operatorname{dist} \left( y , A x + e + N \right) = {\left\Vert \Phi \left( 1 , x , y \right) \right\Vert}_{2} \le \sqrt{3} \operatorname{dist} \left( w \left( x , y \right) , \mathcal{L} \right) .
\end{equation*}


Let $p$ be the point in $A x + e + N$ closest to $y$. Since $y \in D$,
the non-tangency condition gives


\begin{equation*}
\operatorname{dist} \left( p , K^{\star} \left( x \right) \right) \le S^{- 1} \operatorname{dist} \left( p , D \right) \le S^{- 1} {\left\Vert p - y \right\Vert}_{2} ,
\end{equation*}


where the first inequality is immediate if $p \in D$, since then
$p \in K^{\star} \left( x \right)$. The triangle inequality now gives


\begin{equation*}
\operatorname{dist} \left( y , K^{\star} \left( x \right) \right) \le \left( 1 + S^{- 1} \right) \operatorname{dist} \left( y , A x + e + N \right) .
\end{equation*}


Moving $y$ to the nearest point in $K^{\star} \left( x \right)$ changes the reward by
at most ${\left\Vert b \right\Vert}_{2}$ times this distance, which proves the claim.
\end{proof}

We also need to know how close a block average is to the average of the
true conditional means.


\begin{lemma}\label{lem:interval-noise}

  Let $m_{t}$ be the conditional mean of $y_{t}$ given the history and the
  chosen arm on round $t$. For each interval in the first $T$ rounds,
  write $n$ for its length and $\overline{y}$ and $\overline{m}$ for the
  averages of $y_{t}$ and $m_{t}$ over that interval. With probability at
  least $1 - \delta$, the bound



\begin{equation}
n {\left\Vert \overline{y} - \overline{m} \right\Vert}_{2}^{2} \le h_{T}^{2}\label{eq:soft-interval-noise}
\end{equation}



  holds simultaneously on all these intervals.

\end{lemma}


\begin{proof}
Each coordinate of $y_{t} - m_{t}$ is a martingale difference
bounded in absolute value by two. For any fixed interval of length
$n$ and coordinate $i$, Azuma--Hoeffding gives


\begin{equation*}
\operatorname{Pr} \left( \left| {\overline{y}}_{i} - {\overline{m}}_{i} \right| > \frac{h_{T}}{\sqrt{n d_{D}}} \right) \le 2 \operatorname{exp} \left( - \frac{h_{T}^{2}}{8 d_{D}} \right) = \frac{\delta}{d_{D} T^{2}} .
\end{equation*}


There are $d_{D}$ coordinates and at most $T^{2}$ intervals. A union bound
and summing the squared coordinate bounds give the claim.
\end{proof}

We need to know how far a point in $E_{V}$ can be from $\mathcal{L}$.
The next lemma bounds this distance using the noise bound we just
proved. We can then apply Lemma~\ref{lem:subspace-reward-bound} to the outcomes
the learner treats as feasible.


\begin{lemma}\label{lem:subspace-distance-bound}

  With probability at least $1 - \delta$, after any number $J$ of updates
  in Algorithm~\ref{alg:hard-ellipsoid-warmup}, the matrix $V$ satisfies, for every
  $w \in \mathbb{R}^{q}$,


\begin{equation*}
\operatorname{dist} \left( w , \mathcal{L} \right) \le \sqrt{\frac{1 + J h_{T}^{2}}{n}} {\left\Vert w \right\Vert}_{V^{- 1}} .
\end{equation*}


\end{lemma}


\begin{proof}
Number the informative blocks $j = 1 , \ldots , J$. Write $x_{j}$ for
the arm played in block $j$, and ${\overline{y}}_{j}$ and ${\overline{m}}_{j}$
for its average observation and average conditional mean. With
${\overline{w}}_{j} \coloneqq w \left( x_{j} , {\overline{y}}_{j} \right)$, the update rule gives


\begin{equation*}
V = n^{- 1} I_{q} + \sum_{j = 1}^{J} {\overline{w}}_{j} {\overline{w}}_{j}^{T} .
\end{equation*}


Fix $u \in {\mathcal{L}}^{\perp}$ with ${\left\Vert u \right\Vert}_{2} \le 1$. Since the arm is held
fixed within each block, $w \left( x_{j} , {\overline{m}}_{j} \right) \in \mathcal{L}$. Each
informative block has length $n$. By Lemma~\ref{lem:interval-noise}, with
probability at least $1 - \delta$, the bound


\begin{equation*}
{\left( u^{T} {\overline{w}}_{j} \right)}^{2} = {\left( u^{T} \left( 0 , 0 , {\overline{y}}_{j} - {\overline{m}}_{j} \right) \right)}^{2} \le \frac{h_{T}^{2}}{n}
\end{equation*}


holds simultaneously for all informative blocks and all such $u$.
On this event, $u^{T} V u \le \frac{1 + J h_{T}^{2}}{n}$. For any $w$,
Cauchy--Schwarz gives $u^{T} w \le \sqrt{u^{T} V u} {\left\Vert w \right\Vert}_{V^{- 1}}$.
Taking the supremum over the unit ball of ${\mathcal{L}}^{\perp}$ proves the
claim.
\end{proof}


\begin{theorem}\label{thm:hard-ellipsoid-regret}

  Algorithm~\ref{alg:hard-ellipsoid-warmup} satisfies, with probability at least $1 - \delta$,


\begin{equation*}
T V^{\star} - \sum_{t = 1}^{T} r \left( x_{t} , y_{t} \right) \le C_{r} n \left( J_{\text{\normalfont max}} + 1 \right) + \frac{G B_{T} T}{\sqrt{n}} .
\end{equation*}


  Its expected regret is at most the same expression plus
  $C_{r} T \delta$. Taking $n = \left\lceil T^{\frac{2}{3}} \right\rceil$ and
  $\delta = {\left( T + 1 \right)}^{- 2}$ give $R_{T} \le \tilde{O} \left( T^{\frac{2}{3}} \right)$ for fixed problem
  parameters.

\end{theorem}


\begin{proof}
First, there cannot be very many informative blocks. Each one
more than doubles the determinant:


\begin{equation*}
\operatorname{det} \left( V + \overline{w} {\overline{w}}^{T} \right) = \operatorname{det} \left( V \right) \left( 1 + p_{V} \left( x , \overline{y} \right) \right) > 2 \operatorname{det} \left( V \right) .
\end{equation*}


If there are $J$ updates, then $J \le \frac{T}{n}$. Each update adds at most three
to the trace, so $\operatorname{tr} V \le \frac{q}{n} + 3 J \le \frac{q + 3 T}{n}$. The initial determinant
is $n^{- q}$. Using the arithmetic--geometric mean inequality on the
eigenvalues, we get


\begin{equation*}
2^{J} \le \frac{\operatorname{det} V}{\operatorname{det} \left( n^{- 1} I_{q} \right)} \le {\left( 1 + 3 \frac{T}{q} \right)}^{q} , \quad J \le J_{\text{\normalfont max}} .
\end{equation*}


Since $J \le J_{\text{\normalfont max}}$, Lemma~\ref{lem:subspace-distance-bound} gives, with probability
at least $1 - \delta$, simultaneously for every matrix $V$ that the
learner maintains and every $\left( x , y \right) \in X \times D$,



\begin{equation}
\operatorname{dist} \left( w \left( x , y \right) , \mathcal{L} \right) \le \left( \frac{B_{T}}{\sqrt{n}} \right) \sqrt{p_{V} \left( x , y \right)} .\label{eq:hard-subspace-certificate}
\end{equation}



On this event, if $y \in {\hat{K}}_{V} \left( x \right)$, then $p_{V} \left( x , y \right) \le 1$, so
Lemma~\ref{lem:subspace-reward-bound} gives
$r \left( x , y \right) \ge v^{\star} \left( x \right) - G \frac{B_{T}}{\sqrt{n}}$. Taking the minimum over
${\hat{K}}_{V} \left( x \right)$, whenever this set is nonempty, gives



\begin{equation}
H_{V} \left( x \right) \ge v^{\star} \left( x \right) - G \frac{B_{T}}{\sqrt{n}} .\label{eq:hard-optimism}
\end{equation}



Consider a full block in which we do not update the ellipsoid. Its
empirical mean belongs to ${\hat{K}}_{V} \left( x \right)$. Such a block could not have
used the empty-set rule, since that would force the mean outside the
ellipsoid. Thus all the sets ${\hat{K}}_{V} \left( x^{\prime} \right)$ were nonempty, and we
chose $x$ by maximizing $H_{V}$. It follows that


\begin{equation*}
r \left( x , \overline{y} \right) \ge H_{V} \left( x \right) = \operatorname*{max}_{x^{\prime} \in X} H_{V} \left( x^{\prime} \right) \ge V^{\star} - G \frac{B_{T}}{\sqrt{n}} .
\end{equation*}


Since the reward is affine, $r \left( x , \overline{y} \right)$ is the average reward
in the block. The block's regret is therefore at most $G B_{T} \sqrt{n}$.
There are at most $\frac{T}{n}$ such blocks, for a total of $G B_{T} \frac{T}{\sqrt{n}}$.
An informative block may cost as much as $C_{r} n$, but there are at most
$J_{\text{\normalfont max}}$ of them. A final incomplete block also costs at most $C_{r} n$.
Adding these bounds proves the high-probability claim. Outside this
event, regret is still at most $C_{r} T$, so taking expectations
adds at most $C_{r} T \delta$.
\end{proof}


\subsection{Computational tractability and finite-precision implementation}\label{app:hard-computation}

\begin{proof}

We give the implementation promised in Section~\ref{sec:hard-computation}, using Theorem~\ref{thm:bienstock}. The matrix
always satisfies


\begin{equation*}
n^{- 1} I_{q} \preceq V \preceq \left( n^{- 1} + 3 T \right) I_{q} ,
\end{equation*}


since each update adds a positive semidefinite matrix of operator norm
at most three. In particular, $E_{V}$ contains the Euclidean ball of
radius $\frac{1}{\sqrt{n}}$.

Theorem~\ref{thm:bienstock} allows small constraint violations, so its output does not
directly settle the exact emptiness question. For the implementation,
it is enough to find an arm whose original outcome set is empty, or to
certify that every enlarged outcome set is nonempty and plan using those
sets. To leave room for the numerical errors, define

\begin{equation*}
\hat K_V^{\mathrm{in}}(x)\coloneqq\{y\in D:p_V(x,y)\le4\},
\qquad
\hat K_V^{\mathrm{out}}(x)\coloneqq\{y\in D:p_V(x,y)\le9\}.
\end{equation*}

We look for an arm whose inner set is empty, and otherwise plan using
the outer sets. The update rule still uses $p_V>1$.

For a fixed arm $x$, let $S_x\coloneqq\{(1,x,y):y\in D\}$. Its inner
set is empty precisely when $2E_V$ and $S_x$ are disjoint. Since both
sets are compact and convex, this happens if and only if a direction
$z=(z_0,z_X,z_D)$ separates them:

\begin{equation*}
\max_{w\in2E_V}z^T w<\min_{w\in S_x}z^T w.
\end{equation*}

We can compute both sides explicitly. On $S_x$, only $y$ varies, so

\begin{equation*}
\min_{w\in S_x}z^T w
=z_0+z_X^T x+\min_{\|y\|_2\le1}z_D^T y
=z_0+z_X^T x-\|z_D\|_2.
\end{equation*}

For the ellipsoid, write $w=2V^{1/2}u$ with $\|u\|_2\le1$. Then

\begin{equation*}
\max_{w\in2E_V}z^T w
=2\max_{\|u\|_2\le1}(V^{1/2}z)^T u
=2\sqrt{z^T Vz}.
\end{equation*}

Maximizing the separation gap over $\|z\|_2\le1$ gives the Euclidean
distance $d_V(x)\coloneqq\operatorname{dist}(2E_V,S_x)$:

\begin{equation*}
d_V(x)=\max_{\|z\|_2\le1}
\left[z_0+z_X^T x-2\sqrt{z^T Vz}-\|z_D\|_2\right].
\end{equation*}

To justify this identity, write the distance to $2 E_{V}$ as


\begin{equation*}
\operatorname{dist} \left( w , 2 E_{V} \right) = \operatorname*{max}_{{\left\Vert z \right\Vert}_{2} \le 1} \left[ z^{T} w - 2 \sqrt{z^{T} V z} \right] .
\end{equation*}


Minimizing over $w = \left( 1 , x , y \right)$ with $y \in D$ and exchanging the minimum
and maximum gives the claimed expression for $d_{V} \left( x \right)$. The exchange
is valid because the domains are compact and convex, and the
objective is affine in $y$ and concave in $z$.

We maximize this expression jointly over $x\in X$ and $z$, using
Theorem~\ref{thm:bienstock}. After rounding the output to
valid rational points, we obtain an arm $x$ and a certified lower bound
$\ell$ satisfying

\begin{equation*}
\ell\le d_V(x),\qquad
\max_{x'\in X}d_V(x')\le\ell+\eta,
\qquad \eta\coloneqq\frac{1}{10n}.
\end{equation*}

If $\ell>\eta$, then the returned arm has an empty inner set, so we
play it. Otherwise $d_V(x')\le2\eta$ for every arm $x'$. The gap between
$2E_V$ and $3E_V$ then guarantees that every outer set is nonempty, with
a uniform margin for planning. This second outcome does not assert that
every original set is nonempty; nonemptiness of the outer sets is all
we need for the regret argument.

Suppose the empty-set test has not returned an arm. Its additive
guarantee gives $d_{V} \left( x \right) \le 2 \eta$ for every $x$. Thus, for every arm,
there is a $y_{0} \in D$ such that $w_{0} \coloneqq w \left( x , y_{0} \right)$ satisfies


\begin{equation*}
{\left\Vert w_{0} \right\Vert}_{V^{- 1}} \le 2 + 2 \eta \sqrt{n} \le \frac{11}{5} .
\end{equation*}

Define the value used for planning by


\begin{equation*}
H_{V}^{\text{\normalfont out}} \left( x \right) \coloneqq \operatorname*{min}_{y \in {\hat{K}}_{V}^{\text{\normalfont out}} \left( x \right)} r \left( x , y \right) .
\end{equation*}


We will use the following dual formula:


\begin{equation*}
H_{V}^{\text{\normalfont out}} \left( x \right) = a^{T} x + c + \operatorname*{max}_{{\left\Vert z \right\Vert}_{2} \le M} \left[ z_{0} + z_{X}^{T} x - 3 \sqrt{z^{T} V z} - {\left\Vert b + z_{D} \right\Vert}_{2} \right] ,
\end{equation*}


where $M \coloneqq 4 n \left( 1 + {\left\Vert b \right\Vert}_{1} \right)$ suffices. We justify this formula and the bound on $z$ below.

The point $w_{0}$ lies strictly inside $3 E_{V}$. We may therefore apply convex
duality to the constraint $w \left( x , y \right) \in 3 E_{V}$, whose indicator has the
representation


\begin{equation*}
\iota_{3 E_{V}} \left( w \right) = \operatorname*{sup}_{z} \left[ z^{T} w - 3 \sqrt{z^{T} V z} \right] ,
\end{equation*}


where the indicator is zero on the set and $+ \infty$ outside it.
Strong duality holds because $y_{0} \in D$ and $w_{0}$ is in the interior
of $3 E_{V}$ \citep{boyd2004convex}. Minimizing the resulting affine expression
over the unit ball $D$ gives the formula for $H_{V}^{\text{\normalfont out}}$, initially
with a supremum over all $z$.

To bound $z$, write the expression inside that supremum as $F_{x} \left( z \right)$.
Using ${\left\Vert y_{0} \right\Vert}_{2} \le 1$ and Cauchy--Schwarz in the $V$ norm gives


\begin{equation*}
\begin{aligned}
F_{x} \left( z \right) & \le b^{T} y_{0} + z^{T} w_{0} - 3 \sqrt{z^{T} V z} \\
 & \le {\left\Vert b \right\Vert}_{2} - \frac{4}{5} \sqrt{z^{T} V z} \\
 & \le {\left\Vert b \right\Vert}_{2} - \frac{4}{5 \sqrt{n}} {\left\Vert z \right\Vert}_{2} .
\end{aligned}
\end{equation*}


Since $F_{x} \left( 0 \right) = - {\left\Vert b \right\Vert}_{2}$, the supremum is attained at a point with
${\left\Vert z \right\Vert}_{2} \le \left( \frac{5}{2} \right) \sqrt{n} {\left\Vert b \right\Vert}_{2}$. This is smaller than the bound
$M$ used in the implementation.

We maximize the dual formula for $H_{V}^{\text{\normalfont out}}(x)$ jointly over $x\in X$ and $z$ to additive
accuracy $\varepsilon$ to choose the next arm.

We next give the quadratic programs explicitly. For the distance
problem take $\left( \alpha , g , R \right) = \left( 2 , 0 , 1 \right)$, and for the planning problem take
$\left( \alpha , g , R \right) = \left( 3 , b , M \right)$. Replace
$\alpha \sqrt{z^{T} V z} + {\left\Vert g + z_{D} \right\Vert}_{2}$ by $\alpha t + s$ and impose


\begin{equation*}
{\left\Vert x \right\Vert}_{2}^{2} \le 1 , \quad {\left\Vert z \right\Vert}_{2}^{2} \le R^{2} , \quad z^{T} V z \le t^{2} , \quad {\left\Vert g + z_{D} \right\Vert}_{2}^{2} \le s^{2} ,
\end{equation*}



\begin{equation*}
0 \le t \le 1 + \left( 1 + 3 T \right) R , \quad 0 \le s \le 1 + {\left\Vert g \right\Vert}_{1} + R .
\end{equation*}


The remaining objective is $z_{0} + z_{X}^{T} x - \alpha t - s$, with the additional
term $a^{T} x + c$ for planning. The upper bounds leave room beyond the
largest possible values of the two norms. Adding a redundant ball
containing all these bounded variables gives nine quadratic
inequalities, one with a positive-definite quadratic part.
The result in Appendix~\ref{app:qcqp} therefore applies.

A weak solution may slightly violate these constraints. Project its
$x$ and $z$ components onto their balls and round inward to rational
grids. Replace $t$ and $s$ by rational upper approximations to the
corresponding norms. The resulting objective is a valid lower bound
on $d_{V} \left( x \right)$ or $H_{V}^{\text{\normalfont out}} \left( x \right)$ at the returned arm. If the weak
constraint error is $\tau$, projection moves each vector by at most
$\sqrt{\tau}$; the changes in the two norms and the quadratic objective
are bounded by a polynomial in the numerical data and $T$ times
$\sqrt{\tau}$. Thus both the required objective accuracy and a certified
lower bound can be obtained with polynomially many precision bits.
For the distance test, include all these errors in its additive
budget $\eta$: a returned lower bound above $\eta$ then certifies
emptiness, and a lower bound at most $\eta$ implies
$\operatorname*{max}_{x \in X} d_{V} \left( x \right) \le 2 \eta$.

Finally, round each observed outcome inward to a common rational grid,
with Euclidean error at most
$\tau_{y} \coloneqq \frac{{\left( T + 1 \right)}^{- 4}}{1 + {\left\Vert b \right\Vert}_{1}}$. Use these rounded observations in
the block averages and updates. The interval bound in
Lemma~\ref{lem:interval-noise} remains valid with $h_{T}^{2}$ replaced by
$2 h_{T}^{2} + 2 T \tau_{y}^{2}$. The matrix-distance proof and the determinant
argument therefore still apply. An arm chosen by the emptiness test
still forces an update after rounding. An uninformative rounded mean belongs
to $E_{V}$, and hence to both the inner and outer sets. Consequently
such a block must have used the planning rule. Applying
Lemma~\ref{lem:subspace-reward-bound} and the matrix-distance bound on $3 E_{V}$
gives the same regret estimate with the geometric term multiplied
by three, an additional $T \varepsilon$ for planning error, and at most
${\left\Vert b \right\Vert}_{2} T \tau_{y}$ for replacing the actual rewards by rounded ones.

Thus informative and uninformative blocks still cost
$\tilde{O}(n)$ and $\tilde{O}(\sqrt{n})$, respectively. Taking
$\varepsilon=(T+1)^{-2}$ preserves the
$\tilde{O}(T^{2/3})$ regret rate.

The common grids keep all stored matrices rational with polynomial
bit length. Matrix inversion and the comparison $p_{V} > 1$ can then be
performed exactly in polynomial time. There are at most
$\left\lceil \frac{T}{n} \right\rceil$ arm-selection calls, each taking time polynomial in the input
bit length, $T$, and $\operatorname{log} \left( \frac{1}{\varepsilon} \right)$. This proves the claimed
implementation and regret bound.
\end{proof}
}


\clearpage
\draft{
\section{Regret analysis of the $\sqrt{T}$ learner}\label{app:soft-square-root-analysis}

We give the precise statement and proof of Theorem~\ref{thm:soft-square-root}.
We use the coordinates and notation from Section~\ref{sec:soft-square-root},
including $\lambda = \lceil\sqrt{T}\rceil$.

For the analysis, fix $\delta \in \left( 0 , 1 \right)$ and use the quantities
$J_{\text{\normalfont max}} , h_{T} , B_{T} , G , C_{r}$ defined in Appendix~\ref{app:hard-ellipsoid-analysis}.
We use the concentration bound from Lemma~\ref{lem:interval-noise} and adapt the
proof of Lemma~\ref{lem:subspace-distance-bound} to the new matrix update.


\begin{theorem}\label{thm:soft-square-root-regret}

  Consider the variant of Algorithm~\ref{alg:soft-square-root} in which
  each planning call maximizes $U_V$ to within $\varepsilon\ge0$.
  Run in exact arithmetic, it satisfies with probability at least $1-\delta$,


\begin{equation*}
T V^{\star} - \sum_{t = 1}^{T} r \left( x_{t} , y_{t} \right) \le 4 \lambda \left( J_{\text{\normalfont max}} + 1 \right) + \frac{G^{2} B_{T}^{2} T}{4 \lambda} + T \varepsilon ,
\end{equation*}


  With $\varepsilon = \delta = {\left( T + 1 \right)}^{- 2}$, we can
  implement this variant in time polynomial in $\left| \mathcal{M} \right|$ and $T$
  under the weak finite-precision model, with expected regret


\begin{equation*}
R_{T} \le \tilde{O} \left( \left[ q + {\left( 1 + S^{- 1} \right)}^{2} {\left\Vert b \right\Vert}_{2}^{2} q d_{D} \right] \sqrt{T} \right) + O \left( \frac{C_{r}}{T} \right) .
\end{equation*}


  The learner does not need to know $S$ or the analysis quantities
  $J_{\text{\normalfont max}} , h_{T} , B_{T} , G$.

\end{theorem}


\begin{proof}
As before, each completed block at least doubles the
determinant. The stopping rule gives


\begin{equation*}
\operatorname{det} \left( V + n \overline{w} {\overline{w}}^{T} \right) = \operatorname{det} \left( V \right) \left( 1 + n p_{V} \left( x , \overline{y} \right) \right) \ge 2 \operatorname{det} \left( V \right) .
\end{equation*}


Since ${\left\Vert \overline{w} \right\Vert}_{2}^{2} \le 3$, the final trace is at most $q + 3 T$.
Thus, if we complete $J$ blocks,
$2^{J} \le \operatorname{det} V \le {\left( 1 + 3 \frac{T}{q} \right)}^{q}$, so $J \le J_{\text{\normalfont max}}$.

We also need to bound how far $n p_{V} \left( x , {\overline{y}}_{n} \right)$ can jump above
one on the last observation. For $n \ge 2$, convexity gives


\begin{equation*}
n p_{V} \left( x , {\overline{y}}_{n} \right) \le \left( n - 1 \right) p_{V} \left( x , {\overline{y}}_{n - 1} \right) + p_{V} \left( x , y_{n} \right) < 1 + 3 = 4 .
\end{equation*}


The first term is less than one because we did not stop at $n - 1$.
The second is at most three because $V \succeq I_{q}$ and
$p_{V} \left( x , y_{n} \right) \le {\left\Vert w \left( x , y_{n} \right) \right\Vert}_{2}^{2} \le 3$. If we stop at $n = 1$, the bound is
at most three. If the horizon ends before the test succeeds, it is less
than one. Thus every block satisfies



\begin{equation}
n p_{V} \left( x , \overline{y} \right) \le 4 .\label{eq:soft-block-cost}
\end{equation}



We now adapt the proof of Lemma~\ref{lem:subspace-distance-bound}. Number the
completed blocks $j = 1 , \ldots , J$, and let $n_{j}$ be the length of block $j$. Write $x_{j}$
for its arm, ${\overline{y}}_{j}$ and ${\overline{m}}_{j}$ for its average
observation and conditional mean, and
${\overline{w}}_{j} \coloneqq w \left( x_{j} , {\overline{y}}_{j} \right)$. The matrix now has the form


\begin{equation*}
V = I_{q} + \sum_{j = 1}^{J} n_{j} {\overline{w}}_{j} {\overline{w}}_{j}^{T} .
\end{equation*}


By Lemma~\ref{lem:interval-noise}, with probability at least $1 - \delta$, the bound
${\left( u^{T} {\overline{w}}_{j} \right)}^{2} \le \frac{h_{T}^{2}}{n_{j}}$ holds simultaneously for all completed
blocks and all $u \in {\mathcal{L}}^{\perp}$ with ${\left\Vert u \right\Vert}_{2} \le 1$. On this event,


\begin{equation*}
u^{T} V u \le 1 + \sum_{j = 1}^{J} n_{j} \left( \frac{h_{T}^{2}}{n_{j}} \right) = 1 + J h_{T}^{2} \le B_{T}^{2} .
\end{equation*}


This is where the weight $n_{j}$ matters: it cancels the $\frac{1}{n_{j}}$ in the
squared noise bound. Each block contributes at most $h_{T}^{2}$, however
long it is. Applying Cauchy--Schwarz and taking the supremum over these
$u$ gives



\begin{equation}
\operatorname{dist} \left( w , \mathcal{L} \right) \le B_{T} {\left\Vert w \right\Vert}_{V^{- 1}} .\label{eq:soft-subspace-certificate}
\end{equation}



Next, we compare the score with the true value of an arm.
Lemma~\ref{lem:subspace-reward-bound} gives, for every $x \in X$ and $y \in D$,


\begin{equation*}
v^{\star} \left( x \right) \le r \left( x , y \right) + G B_{T} \sqrt{p_{V} \left( x , y \right)} .
\end{equation*}


The error bound grows with $\sqrt{p_{V} \left( x , y \right)}$, while the penalty grows with
$p_{V} \left( x , y \right)$. Writing $u \coloneqq \sqrt{p_{V} \left( x , y \right)}$ and completing the square,


\begin{equation*}
r \left( x , y \right) + \lambda p_{V} \left( x , y \right) \ge v^{\star} \left( x \right) + \lambda u^{2} - G B_{T} u \ge v^{\star} \left( x \right) - \frac{G^{2} B_{T}^{2}}{4 \lambda} .
\end{equation*}


Taking the minimum over $y \in D$, we see that the score underestimates
the true value by at most $\zeta_{T}$:



\begin{equation}
U_{V} \left( x \right) \ge v^{\star} \left( x \right) - \zeta_{T} , \quad \zeta_{T} \coloneqq \frac{G^{2} B_{T}^{2}}{4 \lambda} .\label{eq:soft-optimism}
\end{equation}



We choose an arm whose score is within $\varepsilon$ of the highest score,
so $U_{V} \left( x \right) \ge V^{\star} - \zeta_{T} - \varepsilon$. The block's empirical mean is in
$D$, and the reward is affine. Its regret is therefore at most


\begin{equation*}
n \left( V^{\star} - r \left( x , \overline{y} \right) \right) \le \lambda n p_{V} \left( x , \overline{y} \right) + n \left( \zeta_{T} + \varepsilon \right) \le 4 \lambda + n \left( \zeta_{T} + \varepsilon \right) .
\end{equation*}


Adding this over at most $J_{\text{\normalfont max}} + 1$ blocks proves the high-probability
bound. If the noise bound fails, regret is still at most $C_{r} T$,
so taking expectations adds at most $C_{r} T \delta$.

The difference from the warmup is that we now pay at most $4 \lambda$
in penalty per block, no matter how long the block lasts. We also pay
for the score underestimating the true value, but this costs only
$\zeta_{T}$ per round, which decreases as $\frac{1}{\lambda}$. The two terms to
balance are therefore $\lambda J_{\text{\normalfont max}}$ and $\frac{T}{\lambda}$, giving the
square-root rate.

We still need to implement the planning step. First, rewrite the penalty as


\begin{equation*}
\lambda w^{T} V^{- 1} w = \operatorname*{max}_{z} \left[ z^{T} w - \frac{z^{T} V z}{4 \lambda} \right] .
\end{equation*}


The maximizer is $z = 2 \lambda V^{- 1} w$. Since $V \succeq I_{q}$ and
${\left\Vert w \right\Vert}_{2} \le \sqrt{3}$, we can restrict the maximum to the ball
${\left\Vert z \right\Vert}_{2} \le 4 \lambda$ for every $x \in X$ and $y \in D$. For fixed $x$,
the expression is affine in $y$ and concave in $z$, and both domains
are compact and convex. We can therefore exchange the minimum and
maximum. Writing $z = \left( z_{0} , z_{X} , z_{D} \right)$ gives



\begin{equation}
U_{V} \left( x \right) = a^{T} x + c + \operatorname*{max}_{{\left\Vert z \right\Vert}_{2} \le 4 \lambda} \left[ z_{0} + z_{X}^{T} x - {\left\Vert b + z_{D} \right\Vert}_{2} - \frac{z^{T} V z}{4 \lambda} \right] .\label{eq:soft-conjugate-planner}
\end{equation}




We can now maximize $U_{V}$ by solving the QCQP


\begin{equation*}
\operatorname*{max}_{x , z , s} \quad a^{T} x + c + z_{0} + z_{X}^{T} x - s - \frac{z^{T} V z}{4 \lambda}
\end{equation*}


subject to


\begin{equation*}
{\left\Vert x \right\Vert}_{2}^{2} \le 1 , \quad {\left\Vert z \right\Vert}_{2}^{2} \le 16 \lambda^{2} , \quad {\left\Vert b + z_{D} \right\Vert}_{2}^{2} \le s^{2} , \quad 0 \le s \le H ,
\end{equation*}


where $H \coloneqq 1 + {\left\Vert b \right\Vert}_{1} + 4 \lambda$ is rational. We also add the redundant
constraint


\begin{equation*}
{\left\Vert x \right\Vert}_{2}^{2} + \frac{{\left\Vert z \right\Vert}_{2}^{2}}{16 \lambda^{2}} + \frac{s^{2}}{H^{2}} \le 3 .
\end{equation*}



This gives six quadratic inequalities. The last has a positive-definite
quadratic part in all variables, so we can apply the
weak-optimization result of Bienstock~\citep[Theorem~1.3]{bienstock2016cdt} as described in Section~\ref{app:qcqp}.

The solver may return a point that slightly violates the constraints.
We can make its output $\left( \hat{x} , \hat{z} , \hat{s} \right)$ feasible by projecting
$\hat{x}$ and $\hat{z}$ onto their respective balls. Call the resulting
points $x^{\prime}$ and $z^{\prime}$, and put $s^{\prime} \coloneqq {\left\Vert b + z_{D}^{\prime} \right\Vert}_{2}$. If each constraint
violation is at most $\tau \le 1$, then


\begin{equation*}
{\left\Vert x^{\prime} - \hat{x} \right\Vert}_{2} \le \sqrt{\tau} , \quad {\left\Vert z^{\prime} - \hat{z} \right\Vert}_{2} \le \sqrt{\tau} , \quad s^{\prime} - \hat{s} \le 4 \sqrt{\tau} .
\end{equation*}


The last bound only controls an increase, but this is all we need: the
objective has coefficient $- 1$ on $s$. Since
${\left\Vert V \right\Vert}_{\text{\normalfont op}} \le 1 + 3 T$, the repair decreases the objective by at most
$K_{T} \sqrt{\tau}$, where we can take the rational bound


\begin{equation*}
K_{T} \coloneqq 20 \left( 1 + {\left\Vert a \right\Vert}_{1} + \lambda + T \right) .
\end{equation*}


Choose $\tau \le \operatorname*{min} \left( \frac{1}{4} , {\left( \frac{\varepsilon}{4 \left( K_{T} + 1 \right)} \right)}^{2} \right)$. Including the solver's
objective error, the objective at the repaired point is then within
$\frac{\varepsilon}{2}$ of the optimum. By Equation~\ref{eq:soft-conjugate-planner}, the same lower bound holds for
$U_{V} \left( x^{\prime} \right)$. The number of precision bits we need is polynomial in
$\left| \mathcal{M} \right|$, $\operatorname{log} \left( T + 1 \right)$, and $\operatorname{log} \left( \frac{1}{\varepsilon} \right)$, even when the reward
coefficients are large.

Finally, we round the played arms and observed outcomes to common
rational grids inside their unit balls. Truncating each coordinate
toward zero keeps these points in the balls. The function $U_{V}$ is
Lipschitz in $x$ with constant at most
${\left\Vert a \right\Vert}_{2} + 2 \sqrt{3} \lambda$, so we choose the arm grid finely enough
that rounding costs at most $\frac{\varepsilon}{2}$. For outcomes, keep the
Euclidean error at most $\tau_{y} \coloneqq {\left( T + 1 \right)}^{- 4}$ and use the rounded
observations in the block averages and updates. Their interval averages
still satisfy Equation~\ref{eq:soft-interval-noise} after replacing $h_{T}^{2}$ by
$2 h_{T}^{2} + 2 T \tau_{y}^{2}$. The bounds on the determinant and the penalty per block stay the
same. Replacing actual rewards by rewards at the rounded observations
changes the cumulative reward by at most ${\left\Vert b \right\Vert}_{2} T \tau_{y}$. We
therefore obtain the same expected-regret rate, with universal changes
to the displayed high-probability constants.

The stored matrices are rational, with eigenvalues between $1$ and
$1 + 3 T$. The common grids keep bit lengths polynomial when we form block
averages, update and invert the matrices, and compare the stopping
quantity with its threshold exactly. There are at most $J_{\text{\normalfont max}} + 1$
planning calls, and the QCQP reduction above makes each call run in
polynomial time. This proves the computational claim.
\end{proof}
}


\draft{\section{NP-hardness proofs}\label{app:np-hardness-proofs}}

\draft{\subsection{Planning from learning}\label{app:hardness-planning-from-learning}}



In this section we show that several natural variations of the linear imprecise bandits problem are NP-hard. Thus in some sense, the problem is
the hardest problem that's still tractable.

We first show that efficient planning can be reduced to efficient learning.
Most proofs below show that planning is already NP-hard. Due to the reduction
from planning to learning, this shows that learning is also NP-hard.


\begin{lemma}\label{lem:planning-from-learning}

  {
  Fix a linear imprecise bandit instance $\mathcal{M}$.
  Suppose that, for every hypothesis $z \in Z$, every arm $x$, and every $\eta \in \left( 0 , 1 \right)$, one can compute an outcome
  $y_{z , \eta} \left( x \right)$ satisfying
  }

  {
  $y_{z , \eta} \left( x \right) \in K_{z} \left( x \right)$ and $0 \le r \left( x , y_{z , \eta} \left( x \right) \right) - v_{z} \left( x \right) \le \eta$
  in time polynomial in $| \mathcal{M} |$ and $\frac{1}{\eta}$.
  Suppose also that, for some fixed $\beta > 0$ and polynomial $P$, a
  polynomial-time learner guarantees
  }
  {$R_{T} \le P \left( | \mathcal{M} | \right) T^{1 - \beta}$}

  {
  Then, for every $z \in Z$ and every $\varepsilon \in \left( 0 , 1 \right)$, one can compute in
  time $\operatorname{poly} \left( | \mathcal{M} | , \frac{1}{\varepsilon} \right)$ an arm $\hat{x}$ such
  that with probability at least 2/3, $| V_{z} - v_{z} \left( \hat{x} \right) | \le \varepsilon$.
  }
\end{lemma}


\draft{
\begin{proof}
\draft{
Fix $z \in Z$ and $\varepsilon \in \left( 0 , 1 \right)$, set $\eta \coloneqq \frac{\varepsilon}{6}$, and
simulate the learner on $\mathcal{M}$ for
}

\draft{

\begin{equation*}
T \coloneqq \left\lceil \operatorname*{max} \left( 1 , {\left( 6 \frac{P \left( | \mathcal{M} | \right)}{\varepsilon} \right)}^{\frac{1}{\beta}} \right) \right\rceil
\end{equation*}

}

\draft{
rounds. Whenever the learner plays $x_{t}$, let nature use the point mass at
$y_{z , \eta} \left( x_{t} \right)$. This is allowed because
$y_{z , \eta} \left( x_{t} \right) \in K_{z} \left( x_{t} \right)$. Record the rewards and choose the round
with the largest reward:
}

\draft{

\begin{equation*}
{\tilde{v}}_{t} \coloneqq r \left( x_{t} , y_{z , \eta} \left( x_{t} \right) \right) , \quad \hat{t} \in \operatorname{argmax}_{t = 1 , \ldots , T} {\tilde{v}}_{t} ,
\end{equation*}

}

\draft{
and return $\hat{x} \coloneqq x_{\hat{t}}$.
}

\draft{
The approximation guarantee gives
$v_{z} \left( \hat{x} \right) \ge {\tilde{v}}_{\hat{t}} - \eta$. Our choice of $\hat{t}$ therefore gives
}

\draft{

\begin{equation*}
0 \le V_{z} - v_{z} \left( \hat{x} \right) \le V_{z} - {\tilde{v}}_{\hat{t}} + \eta \le V_{z} - \frac{1}{T} \sum_{t = 1}^{T} {\tilde{v}}_{t} + \eta .
\end{equation*}

}

\draft{
Taking expectations and using the definition of $R_{T}$ for this nature policy,
}

\draft{

\begin{equation*}
\mathbb{E} \left[ V_{z} - v_{z} \left( \hat{x} \right) \right] \le \frac{R_{T}}{T} + \eta \le P \left( | \mathcal{M} | \right) T^{- \beta} + \eta \le \frac{\varepsilon}{3} .
\end{equation*}

}

\draft{
The gap $V_{z} - v_{z} \left( \hat{x} \right)$ is nonnegative, so Markov's inequality shows
that it exceeds $\varepsilon$ with probability at most $\frac{1}{3}$.
The horizon and the simulation time are polynomial in $| \mathcal{M} |$ and $\frac{1}{\varepsilon}$
because $\beta$ is fixed.
}
\end{proof}
}

Below we show that the ``smoothness'' of $X$ and $D$ are crucial for efficiency. It is easy to land in the NP-hard territory if we make either $X$ or $D$ a convex shape with edges and corners.





\draft{\subsection{Proof of Theorem~\ref{thm:additive-simplex-np-hardness}}\label{app:hardness-simplex}}



First we show that if, instead of a Euclidean ball, we allow $D$ to be a simplex, then even known hypothesis planning becomes NP-hard. Then, using the planning to learning reduction (Lemma~\ref{lem:planning-from-learning}), we conclude that learning is NP-hard.







\begin{proof}
We reduce from MAX-CUT, as defined in Appendix~\ref{app:max-cut}. Let $\left( G , k \right)$ be an
instance, where $G = \left( V , E \right)$ has $n$ vertices and $m \ge 1$ edges and
$k \in \left\{ 1 , \ldots , m \right\}$. Orient the edges arbitrarily and let
$B_{G} \in \mathbb{R}^{m \times n}$ be the resulting edge-vertex incidence matrix, i.e., if
$e = \left\{ i , j \right\}$ is oriented from $j$ to $i$, then row $e$ has $1$ in column $i$,
$- 1$ in column $j$, and zeros elsewhere. Put


\begin{equation*}
X \coloneqq \left\{ x \in \mathbb{R}^{m} : {\left\Vert x \right\Vert}_{2} \le 1 \right\} , \quad Z \coloneqq \left[ 1 , 2 \right] .
\end{equation*}


Write an outcome as $y = \left( p , q , s \right) \in \mathbb{R}^{n} \times \mathbb{R}^{n} \times \mathbb{R}$ and take


\begin{equation*}
D \coloneqq \left\{ \left( p , q , s \right) : p_{i} \ge 0 , q_{i} \ge 0 \text{\normalfont  for } i = 1 , \ldots , n , s \ge 0 , \sum_{i = 1}^{n} \left( p_{i} + q_{i} \right) + s = 1 \right\} .
\end{equation*}


Let $\gamma > 0$ be a constant to be chosen later. For every $z \in Z$, let



\begin{equation*}
K_{z} \left( x \right) \coloneqq \left\{ \left( p , q , s \right) \in D : p - q = \gamma B_{G}^{T} x \right\} .
\end{equation*}


This is of the form in Equation~\ref{eq:setting-compatible-set}, where


\begin{equation*}
C_{z} \left( p , q , s \right) \coloneqq p - q , \quad B_{z} x \coloneqq - \gamma B_{G}^{T} x , \quad d_{z} \coloneqq 0 .
\end{equation*}


Let the reward be


\begin{equation*}
r \left( x , \left( p , q , s \right) \right) \coloneqq \sum_{i = 1}^{n} \left( p_{i} + q_{i} \right) .
\end{equation*}


The maps $z \mapsto B_{z}$, $z \mapsto C_{z}$, and $z \mapsto d_{z}$ are constant,
hence affine, as required, and the reward lies in $\left[ 0 , 1 \right]$ on $D$.

In particular, $K_{z} \left( x \right)$ is independent of $z$.

Next we show that, for every $z \in Z$, $\left( G , k \right)$ is a yes-instance of MAX-CUT
if and only if the constructed imprecise bandit instance has an arm $x \in X$
satisfying



\begin{equation*}
v_{z} \left( x \right) \ge 2 \gamma \sqrt{k} .
\end{equation*}


A vector $\sigma \in {\left\{ - 1 , 1 \right\}}^{n}$ represents a partition of the vertices of $G$
into the two sides $\left\{ i : \sigma_{i} = 1 \right\}$ and $\left\{ i : \sigma_{i} = - 1 \right\}$. Let
$c_{G} \left( \sigma \right)$ denote the number of edges between these sides.


\begin{lemma}\label{lem:simplex-cut-norm}

  For every $\sigma \in {\left\{ - 1 , 1 \right\}}^{n}$,


\begin{equation*}
{\left\Vert B_{G} \sigma \right\Vert}_{2} = 2 \sqrt{c_{G} \left( \sigma \right)} .
\end{equation*}


\end{lemma}


\begin{proof}
Fix $\sigma \in {\left\{ - 1 , 1 \right\}}^{n}$. For every edge $e = \left\{ i , j \right\}$,


\begin{equation*}
{\left( B_{G} \sigma \right)}_{e}^{2} = {\left( \sigma_{i} - \sigma_{j} \right)}^{2} = \begin{cases}
4 & \text{\normalfont if } \sigma_{i} \ne \sigma_{j} \\
0 & \text{\normalfont otherwise}
\end{cases} .
\end{equation*}


Summing over $e \in E$ and taking the square root proves the identity.
\end{proof}

For the forward implication, suppose $\left( G , k \right)$ is a yes-instance. Choose
$\sigma$ such that $c_{G} \left( \sigma \right) \ge k$, and define the arm


\begin{equation*}
x_{\sigma} \coloneqq B_{G} \frac{\sigma}{{\left\Vert B_{G} \sigma \right\Vert}_{2}} .
\end{equation*}


Since $c_{G} \left( \sigma \right) \ge k \ge 1$, Lemma~\ref{lem:simplex-cut-norm} shows that $x_{\sigma}$ is
well-defined. Moreover, we get the following inequality.


\begin{lemma}\label{lem:simplex-duality}

  For every $\sigma \in {\left\{ - 1 , 1 \right\}}^{n}$ such that $B_{G} \sigma \ne 0$, we have
  ${\left\Vert B_{G}^{T} x_{\sigma} \right\Vert}_{1} \ge {\left\Vert B_{G} \sigma \right\Vert}_{2}$. Moreover, equality holds when
  $\sigma$ maximizes ${\left\Vert B_{G} \sigma \right\Vert}_{2}$ over ${\left\{ - 1 , 1 \right\}}^{n}$; in particular,
  equality holds for at least one $\sigma$.

\end{lemma}


\draft{
\begin{proof}
\draft{
Choosing the sign of each coordinate gives
}

\draft{

\begin{equation*}
{\left\Vert B_{G}^{T} x_{\sigma} \right\Vert}_{1} = \operatorname*{max}_{\tau \in {\left\{ - 1 , 1 \right\}}^{n}} \tau^{T} B_{G}^{T} x_{\sigma} = \operatorname*{max}_{\tau \in {\left\{ - 1 , 1 \right\}}^{n}} x_{\sigma}^{T} B_{G} \tau .
\end{equation*}

}

\draft{
Taking $\tau = \sigma$ yields
${\left\Vert B_{G}^{T} x_{\sigma} \right\Vert}_{1} \ge x_{\sigma}^{T} B_{G} \sigma = {\left\Vert B_{G} \sigma \right\Vert}_{2}$.
Now let $\sigma^{\star}$ maximize ${\left\Vert B_{G} \sigma \right\Vert}_{2}$. Since $m \ge 1$, this maximum
is positive, so $x_{\sigma^{\star}}$ is well-defined. For every
$\tau \in {\left\{ - 1 , 1 \right\}}^{n}$, Cauchy--Schwarz gives
}

\draft{

\begin{equation*}
x_{\sigma^{\star}}^{T} B_{G} \tau \le {\left\Vert x_{\sigma^{\star}} \right\Vert}_{2} {\left\Vert B_{G} \tau \right\Vert}_{2} \le {\left\Vert B_{G} \sigma^{\star} \right\Vert}_{2} .
\end{equation*}

}

\draft{
Taking the maximum over $\tau$ gives the reverse inequality, proving
equality.
}
\end{proof}
}

The lemma above gives us a way to link properties of cuts with properties of arms. The missing ingredient is to show that the value of an arm has something to do with ${\left\Vert B_{G}^{T} x \right\Vert}_{1}$. We show this in the next two lemmas.

By construction, $x_{\sigma} \in X$. Since $K_{z} \left( x \right)$ is independent of $z$, write
$K \left( x \right) \coloneqq K_{z} \left( x \right) = \left\{ \left( p , q , s \right) \in D : p - q = \gamma B_{G}^{T} x \right\} .$


\begin{lemma}

  For every $x \in X$, the vector $t \coloneqq \gamma B_{G}^{T} x$ satisfies
  ${\left\Vert t \right\Vert}_{1} \le \gamma \sqrt{2 m n}$. Therefore, if
  $\gamma \le \frac{1}{\sqrt{2 m n}}$, then ${\left\Vert t \right\Vert}_{1} \le 1$.

\end{lemma}


\draft{
\begin{proof}
\draft{
Each edge coordinate $x_{e}$ contributes $x_{e}$ and $- x_{e}$ to two
coordinates of $B_{G}^{T} x$ before contributions from different edges are
added. The triangle inequality therefore gives
}

\draft{

\begin{equation*}
{\left\Vert t \right\Vert}_{1} = \gamma {\left\Vert B_{G}^{T} x \right\Vert}_{1} \le 2 \gamma {\left\Vert x \right\Vert}_{1} \le 2 \gamma \sqrt{m} {\left\Vert x \right\Vert}_{2} \le 2 \gamma \sqrt{m} \le \gamma \sqrt{2 m n} ,
\end{equation*}

}

\draft{
where the second inequality is Cauchy--Schwarz, the third uses
${\left\Vert x \right\Vert}_{2} \le 1$, and the last uses $n \ge 2$, which follows from $m \ge 1$.
}
\end{proof}
}


\begin{lemma}\label{lem:simplex-arm-value}

  Suppose $\gamma \le \frac{1}{\sqrt{2 m n}}$. For each arm $x \in X$, set
  $t \coloneqq \gamma B_{G}^{T} x$. Then $K \left( x \right)$ is nonempty, and for all $z$,
  $v_{z} \left( x \right) = \operatorname*{min}_{y \in K \left( x \right)} r \left( x , y \right) = {\left\Vert t \right\Vert}_{1} = \gamma {\left\Vert B_{G}^{T} x \right\Vert}_{1}$. Moreover, the minimum is attained at the outcome $y = \left( p , q , s \right)$ given by
  $p_{i} \coloneqq \operatorname*{max} \left( t_{i} , 0 \right)$, $q_{i} \coloneqq \operatorname*{max} \left( - t_{i} , 0 \right)$, and $s \coloneqq 1 - {\left\Vert t \right\Vert}_{1}$.

\end{lemma}


\draft{
\begin{proof}
\draft{
Fix $x \in X$ and set $t \coloneqq \gamma B_{G}^{T} x$. The preceding lemma
gives ${\left\Vert t \right\Vert}_{1} \le 1$. For the outcome in the statement, $p , q , s$ are
nonnegative, $p - q = t$, and
}

\draft{

\begin{equation*}
\sum_{i = 1}^{n} \left( p_{i} + q_{i} \right) + s = {\left\Vert t \right\Vert}_{1} + 1 - {\left\Vert t \right\Vert}_{1} = 1 .
\end{equation*}

}

\draft{
Thus this outcome belongs to $K \left( x \right)$. For any $\left( p , q , s \right) \in K \left( x \right)$,
nonnegativity gives $p_{i} + q_{i} \ge \left| p_{i} - q_{i} \right| = \left| t_{i} \right|$ for every $i$, and hence
$r \left( x , \left( p , q , s \right) \right) \ge {\left\Vert t \right\Vert}_{1}$. The outcome in the statement attains equality in
every coordinate, so its reward is the minimum.
}
\end{proof}
}

Note that this also shows that given an arm $x$, we can compute $v_{z} \left( x \right)$ efficiently, which is one of the prerequisites for the application of Lemma~\ref{lem:planning-from-learning}.

For the final step, apply Lemma~\ref{lem:simplex-arm-value}, Lemma~\ref{lem:simplex-duality}, and
Lemma~\ref{lem:simplex-cut-norm} to obtain


\begin{equation*}
v_{z} \left( x_{\sigma} \right) = \gamma {\left\Vert B_{G}^{T} x_{\sigma} \right\Vert}_{1} \ge \gamma {\left\Vert B_{G} \sigma \right\Vert}_{2} = 2 \gamma \sqrt{c_{G} \left( \sigma \right)} \ge 2 \gamma \sqrt{k} .
\end{equation*}


For the reverse implication, suppose there is an arm $x \in X$ satisfying
$v_{z} \left( x \right) \ge 2 \gamma \sqrt{k}$. By $\ell_{1}$--$\ell_{\infty}$ duality, choose
$\sigma \in {\left\{ - 1 , 1 \right\}}^{n}$ such that


\begin{equation*}
{\left\Vert B_{G}^{T} x \right\Vert}_{1} = \sigma^{T} B_{G}^{T} x .
\end{equation*}


Using Lemma~\ref{lem:simplex-arm-value} and Lemma~\ref{lem:simplex-cut-norm},


\begin{equation*}
2 \gamma \sqrt{k} \le v_{z} \left( x \right) = \gamma {\left\Vert B_{G}^{T} x \right\Vert}_{1} = \gamma x^{T} B_{G} \sigma \le \gamma {\left\Vert x \right\Vert}_{2} {\left\Vert B_{G} \sigma \right\Vert}_{2} \le 2 \gamma \sqrt{c_{G} \left( \sigma \right)} .
\end{equation*}


Since $\gamma > 0$, this implies $c_{G} \left( \sigma \right) \ge k$, so $\left( G , k \right)$ is a yes-instance.
This proves the claimed equivalence.

To apply Lemma~\ref{lem:planning-from-learning}, we need to show that \emph{approximate}
planning is hard. Although the lemma returns only an approximate optimizer,
this is sufficient because the maximum cut size $M_{G}$ is an integer between
$0$ and $m$, and the corresponding planning values are
$V_{z} = 2 \gamma \sqrt{M_{G}}$. Taking $\gamma \coloneqq \frac{1}{2 m n}$ ensures that every
$K_{z} \left( x \right)$ is nonempty and that consecutive planning values are separated by
at least $\frac{1}{2 m^{2} n}$. Thus an additive approximation with error
$\varepsilon \coloneqq \frac{1}{6 m^{2} n}$ distinguishes $M_{G} \ge k$ from $M_{G} \le k - 1$.

By Lemma~\ref{lem:simplex-arm-value}, we can compute both the value of the returned arm
and a minimizing outcome in polynomial time. We can therefore decide MAX-CUT
by comparing the returned arm's value with the midpoint between
$2 \gamma \sqrt{k - 1}$ and $2 \gamma \sqrt{k}$, computed to sufficient precision.
Hence Lemma~\ref{lem:planning-from-learning} implies that a randomized polynomial-time
learner would place MAX-CUT in $\text{\normalfont BPP}$, and therefore that
$\text{\normalfont NP} \subseteq \text{\normalfont BPP}$. A deterministic learner would instead imply
$\text{\normalfont P} = \text{\normalfont NP}$.
\end{proof}




\draft{\subsection{Proof of Theorem~\ref{thm:additive-polytope-np-hardness}}\label{app:hardness-polytope}}



\draft{
We now keep the Euclidean outcome ball and the additive bilinear constraint,
but allow the arm set to be a polytope given by rational inequalities.
Planning is NP-hard even when the hypothesis is known and the arm set is
the cube. The cube has only $2 n$ defining inequalities, but $2^{n}$ vertices.
}







\draft{
\begin{proof}
\draft{
We again reduce from MAX-CUT. Let $G = \left( V , E \right)$ have $n$ vertices and
$m \ge 1$ edges. Orient the edges arbitrarily and let
$B_{G} \in \mathbb{R}^{m \times n}$ be the edge-vertex incidence matrix, so that
}

\draft{

\begin{equation*}
{\left( B_{G} x \right)}_{e} = x_{i} - x_{j}
\end{equation*}

}

\draft{
for an edge $e = \left\{ i , j \right\}$. Put
}

\draft{

\begin{equation*}
\varepsilon \coloneqq \frac{1}{4 m} , \quad X \coloneqq {\left[ - 1 , 1 \right]}^{n} , \quad Z \coloneqq \left[ 1 , 2 \right] .
\end{equation*}

}

\draft{
Write outcomes as $y = \left( u , s \right) \in \mathbb{R}^{m} \times \mathbb{R}$ and take
}

\draft{

\begin{equation*}
D \coloneqq \left\{ \left( u , s \right) : {\left\Vert u \right\Vert}_{2}^{2} + s^{2} \le 1 \right\} .
\end{equation*}

}

\draft{
With $W = \mathbb{R}^{m}$, define
}

\draft{

\begin{equation*}
F_{0} \left( x , z \right) \coloneqq - \varepsilon z B_{G} x , \quad F_{1} \left( \left( u , s \right) , z \right) \coloneqq z u ,
\end{equation*}

}

\draft{
The constraint is $F_{0} \left( x , z \right) + F_{1} \left( \left( u , s \right) , z \right) = 0$, and the reward is
$r \left( x , \left( u , s \right) \right) \coloneqq s$. Both constraint maps are bilinear and the reward is
linear. Since $z \ne 0$, the constraint is equivalent to
}

\draft{

\begin{equation*}
u = \varepsilon B_{G} x ,
\end{equation*}

}

\draft{
so the feasible outcomes do not depend on $z$. Every $x \in X$ satisfies
}

\draft{

\begin{equation*}
{\left\Vert B_{G} x \right\Vert}_{2}^{2} = \sum_{\left\{ i , j \right\} \in E} {\left( x_{i} - x_{j} \right)}^{2} \le 4 m ,
\end{equation*}

}

\draft{
so
}

\draft{

\begin{equation*}
{\left\Vert \varepsilon B_{G} x \right\Vert}_{2}^{2} \le \frac{1}{4 m} \le \frac{1}{4} .
\end{equation*}

}

\draft{
Thus the feasible outcome set $K_{z} \left( x \right)$ is nonempty for every $x \in X$
and $z \in Z$. Minimizing $s$ over this set gives
}

\draft{

\begin{equation*}
v_{z} \left( x \right) = - \sqrt{1 - \varepsilon^{2} {\left\Vert B_{G} x \right\Vert}_{2}^{2}} .
\end{equation*}

}

\draft{
Set $q \left( x \right) \coloneqq {\left\Vert B_{G} x \right\Vert}_{2}^{2}$. With every coordinate except $x_{i}$ fixed,
$q \left( x \right)$ is a convex quadratic in $x_{i}$, so replacing $x_{i}$ by one of the
endpoints of $\left[ - 1 , 1 \right]$ does not decrease $q$. Applying this replacement one
coordinate at a time produces a sign vector $\sigma \in {\left\{ - 1 , 1 \right\}}^{n}$ with
$q \left( \sigma \right) \ge q \left( x \right)$. At a sign vector,
}

\draft{

\begin{equation*}
q \left( \sigma \right) = \sum_{\left\{ i , j \right\} \in E} {\left( \sigma_{i} - \sigma_{j} \right)}^{2} = 4 \, \left| \text{\normalfont cut} \left( \sigma \right) \right| .
\end{equation*}

}

\draft{
Consequently,
}

\draft{

\begin{equation*}
\operatorname*{max}_{x \in {\left[ - 1 , 1 \right]}^{n}} q \left( x \right) = 4 \text{\normalfont MaxCut} \left( G \right) .
\end{equation*}

}

\draft{
Since $v_{z} \left( x \right)$ is strictly increasing in $q \left( x \right)$, if the maximum cut has
size $k$, the optimal robust value is
}

\draft{

\begin{equation*}
V_{k} = - \sqrt{1 - \frac{k}{4 m^{2}}} .
\end{equation*}

}

\draft{
For $k = 1 , \ldots , m$, consecutive possible values satisfy
}

\draft{

\begin{equation*}
V_{k} - V_{k - 1} = \frac{\frac{1}{4 m^{2}}}{\sqrt{1 - \frac{k - 1}{4 m^{2}}} + \sqrt{1 - \frac{k}{4 m^{2}}}} \ge \frac{1}{8 m^{2}} .
\end{equation*}

}

\draft{
Hence an estimate of $V_{z}$ within
}

\draft{

\begin{equation*}
\Delta \coloneqq \frac{1}{32 m^{2}}
\end{equation*}

}

\draft{
identifies $k$ after we compute the $m + 1$ candidate values to polynomially
many bits. Thus approximating the optimal robust value is NP-hard.
}

\draft{
We now show how to use the learner to find a maximum cut. To simulate
the learner, we need rational outcomes in $K_{z} \left( x \right)$. The outcome attaining
$v_{z} \left( x \right)$ need not be rational, so we compute a rational feasible
approximation. Suppose the learner in the theorem exists and set
}

\draft{

\begin{equation*}
\eta \coloneqq \frac{\Delta}{6} , \quad T \coloneqq \left\lceil {\left( 6 \frac{P \left( | \mathcal{M} | \right)}{\Delta} \right)}^{\frac{1}{\beta}} \right\rceil .
\end{equation*}

}

\draft{
In the weak bit model the learner outputs rationally encoded arms. For each
played arm $x_{t}$, we can use binary search to compute a rational $a_{t}$ in
polynomial time such that
}

\draft{

\begin{equation*}
0 \le \sqrt{1 - \varepsilon^{2} q \left( x_{t} \right)} - a_{t} \le \eta , \quad a_{t}^{2} \le 1 - \varepsilon^{2} q \left( x_{t} \right) .
\end{equation*}

}

\draft{
Let nature return the point
}

\draft{

\begin{equation*}
y_{t} \coloneqq \left( \varepsilon B_{G} x_{t} , - a_{t} \right) .
\end{equation*}

}

\draft{
This is a rational point of $K_{z} \left( x_{t} \right)$, and its reward satisfies
}

\draft{

\begin{equation*}
0 \le r \left( x_{t} , y_{t} \right) - v_{z} \left( x_{t} \right) \le \eta .
\end{equation*}

}

\draft{
Writing $V_{z} = V_{k}$ and
}

\draft{

\begin{equation*}
S_{T} \coloneqq \sum_{t = 1}^{T} \left( V_{z} - v_{z} \left( x_{t} \right) \right) ,
\end{equation*}

}

\draft{
the definition of expected regret gives
}

\draft{

\begin{equation*}
E \left[ S_{T} \right] \le R_{T} + T \eta .
\end{equation*}

}

\draft{
Applying the learner's regret guarantee,
}

\draft{

\begin{equation*}
E \left[ S_{T} \right] \le P \left( | \mathcal{M} | \right) T^{1 - \beta} + T \eta .
\end{equation*}

}

\draft{
If the learner never plays a $\Delta$-optimal arm, then $S_{T} > \Delta T$.
Markov's inequality and the choices of $T$ and $\eta$ imply
}

\draft{

\begin{equation*}
\operatorname{Pr} \left( \text{\normalfont no } \Delta \text{\normalfont -optimal arm is played} \right) \le \frac{P \left( | \mathcal{M} | \right)}{\Delta T^{\beta}} + \frac{\eta}{\Delta} \le \frac{1}{3} .
\end{equation*}

}

\draft{
Select a played arm maximizing the rational quantity $q \left( x_{t} \right)$ and round its
coordinates to endpoints without decreasing $q$. With probability at least
$\frac{2}{3}$, the resulting sign vector has robust value at least $V_{k} - \Delta$. A cut
of size at most $k - 1$ would instead have robust value at most
$V_{k - 1} \le V_{k} - \frac{1}{8 m^{2}} < V_{k} - \Delta$. The rounded sign vector therefore encodes
a maximum cut. The horizon, the square-root approximations, and the rounding
all have polynomial bit complexity, proving $\text{\normalfont NP} \subseteq \text{\normalfont BPP}$. If the
learner is deterministic, the same argument succeeds with certainty and gives
$\text{\normalfont P} = \text{\normalfont NP}$.
}

\draft{
It remains to verify non-tangency. For a fixed $x$, put
}

\draft{

\begin{equation*}
a \coloneqq {\left\Vert \varepsilon B_{G} x \right\Vert}_{2} , \quad \rho \coloneqq \sqrt{1 - a^{2}} \ge \frac{\sqrt{3}}{2} .
\end{equation*}

}

\draft{
The outcomes satisfying the affine constraint form the line $L_{z} \left( x \right)$,
whose intersection with $D$ is $K_{z} \left( x \right)$:
}

\draft{

\begin{equation*}
L_{z} \left( x \right) = \left\{ \left( \varepsilon B_{G} x , s \right) : s \in \mathbb{R} \right\} , \quad K_{z} \left( x \right) = \left\{ \left( \varepsilon B_{G} x , s \right) : \left| s \right| \le \rho \right\} .
\end{equation*}

}

\draft{
If $p = \left( \varepsilon B_{G} x , s \right) \in L_{z} \left( x \right)$ lies outside $D$ and $t \coloneqq \left| s \right| > \rho$,
then
}

\draft{

\begin{equation*}
\operatorname{dist} \left( p , K_{z} \left( x \right) \right) = t - \rho , \quad \operatorname{dist} \left( p , D \right) = \sqrt{a^{2} + t^{2}} - 1 .
\end{equation*}

}

\draft{
Factoring the second expression gives
}

\draft{

\begin{equation*}
\frac{\operatorname{dist} \left( p , D \right)}{\operatorname{dist} \left( p , K_{z} \left( x \right) \right)} = \frac{t + \rho}{\sqrt{a^{2} + t^{2}} + 1} \ge \rho \ge \frac{\sqrt{3}}{2} .
\end{equation*}

}

\draft{
To check the first inequality, multiply by the positive denominator,
subtract $\rho$ from both sides, and square. This gives the equivalent inequality
$\left( 1 - \rho^{2} \right) \left( t^{2} - \rho^{2} \right) \ge 0$. Thus the uniform non-tangency condition holds with
$S = \frac{\sqrt{3}}{2}$.
}
\end{proof}
}



\draft{\subsection{Proof of Theorem~\ref{thm:trilinear-balls-np-hardness}}\label{app:hardness-trilinear}}



\draft{
We keep the Euclidean arm and outcome balls, but allow the coefficients in
the equations for $y$ to depend on $x$. We add an outcome coordinate
that is fixed to one, so that the constraints are linear in the full
outcome vector with no constant term. Let $Y$ be the ambient outcome
space, let $W$ be the constraint space, and let
}

\draft{

\begin{equation*}
F : \mathbb{R}^{d_{X}} \times \mathbb{R}^{d_{Z}} \times Y \to W
\end{equation*}

}

\draft{
be linear in each argument separately. For each $x$ and $z$, write
$F_{x , z} \left( y \right) \coloneqq F \left( x , z , y \right)$ and define
}

\draft{

\begin{equation*}
L_{z} \left( x \right) \coloneqq \operatorname{ker} F_{x , z} , \quad K_{z} \left( x \right) \coloneqq L_{z} \left( x \right) \cap D , \quad v_{z} \left( x \right) \coloneqq \operatorname*{min}_{y \in K_{z} \left( x \right)} r \left( x , y \right) , \quad V_{z} \coloneqq \operatorname*{max}_{x \in X} v_{z} \left( x \right) .
\end{equation*}

}

\draft{
Here $D$ is a Euclidean ball in the hyperplane where the added
coordinate equals one. We will show that planning is now NP-hard even
when $z$ is known.
}

\draft{
Let $G = \left( V , E \right)$ be an undirected graph with $V = \left\{ 1 , \ldots , n \right\}$ and
$E = \left\{ e_{1} , \ldots , e_{m} \right\}$, where $m \ge 1$ and $e_{\ell} = \left\{ i_{\ell} , j_{\ell} \right\}$. Put $d \coloneqq n + m$.
For $q = \left( u , w \right) \in \mathbb{R}^{n} \times \mathbb{R}^{m}$, define the homogeneous cubic
}

\draft{

\begin{equation*}
p_{G} \left( q \right) \coloneqq \sum_{\ell = 1}^{m} u_{i_{\ell}} u_{j_{\ell}} w_{\ell} .
\end{equation*}

}


\begin{lemma}\label{lem:clique-cubic}

  \draft{
  If $\omega \left( G \right)$ is the clique number of $G$, then
  }

  \draft{

\begin{equation*}
\operatorname*{max}_{{\left\Vert q \right\Vert}_{2} = 1} p_{G} {\left( q \right)}^{2} = \frac{2}{27} \left( 1 - \frac{1}{\omega \left( G \right)} \right) .
\end{equation*}

  }

\end{lemma}


\draft{
\begin{proof}
\draft{
Write $\rho \coloneqq {\left\Vert u \right\Vert}_{2}$ and $\sigma \coloneqq {\left\Vert w \right\Vert}_{2}$. For fixed $u$ and
$\sigma$, Cauchy--Schwarz gives
}

\draft{

\begin{equation*}
\operatorname*{max}_{{\left\Vert w \right\Vert}_{2} = \sigma} p_{G} \left( u , w \right) = \sigma \sqrt{\sum_{\ell = 1}^{m} u_{i_{\ell}}^{2} u_{j_{\ell}}^{2}} .
\end{equation*}

}

\draft{
Write $u = \rho s$ with ${\left\Vert s \right\Vert}_{2} = 1$ and set $\pi_{i} \coloneqq s_{i}^{2}$. The
Motzkin--Straus theorem \citep{motzkin1965maxima} gives
}

\draft{

\begin{equation*}
\operatorname*{max}_{{\left\Vert s \right\Vert}_{2} = 1} \sum_{\left\{ i , j \right\} \in E} s_{i}^{2} s_{j}^{2} = \operatorname*{max}_{\pi_{i} \ge 0 , \sum_{i} \pi_{i} = 1} \sum_{\left\{ i , j \right\} \in E} \pi_{i} \pi_{j} = \frac{1}{2} \left( 1 - \frac{1}{\omega \left( G \right)} \right) .
\end{equation*}

}

\draft{
Finally,
}

\draft{

\begin{equation*}
\operatorname*{max}_{\rho^{2} + \sigma^{2} = 1 , \rho \ge 0 , \sigma \ge 0} \rho^{2} \sigma = \frac{2}{3 \sqrt{3}} .
\end{equation*}

}

\draft{
Substituting these bounds into the first display and squaring proves
the claim. Since $p_{G}$ is odd, its
maximum equals its maximum absolute value.
}
\end{proof}
}







\draft{
\begingroup
\renewcommand{\proofname}{Proof of Theorem~\ref{thm:trilinear-balls-np-hardness}}
\begin{proof}
\draft{
We may restrict CLIQUE to instances with $m \ge 1$ and
$3 \le k \le n$, since we can decide the other cases in polynomial time. Given
$\left( G , k \right)$, use the cubic $p_{G}$ above and write an arm as
$x = \left( x_{0} , \xi \right) \in \mathbb{R} \times \mathbb{R}^{d}$. Take
}

\draft{

\begin{equation*}
X \coloneqq \left\{ \left( x_{0} , \xi \right) : {\left( x_{0} - \frac{5}{3} \right)}^{2} + {\left\Vert \xi \right\Vert}_{2}^{2} \le 1 \right\} .
\end{equation*}

}

\draft{
This is a full-dimensional Euclidean ball and $x_{0} \ge \frac{2}{3}$ throughout $X$.
The ratio $q \coloneqq \frac{\xi}{x_{0}}$ ranges over exactly the ball
${\left\Vert q \right\Vert}_{2} \le \frac{3}{4}$. Indeed,
}

\draft{

\begin{equation*}
1 - {\left( x_{0} - \frac{5}{3} \right)}^{2} - \frac{9}{16} x_{0}^{2} = - \frac{{\left( 15 x_{0} - 16 \right)}^{2}}{144} \le 0 ,
\end{equation*}

}

\draft{
so every ratio has norm at most $\frac{3}{4}$. Conversely, for any such
$q$, the choice $x_{0} = \frac{16}{15}$ and $\xi = x_{0} q$ belongs to $X$.
}

\draft{
Let $Z \coloneqq \left[ 1 , 2 \right]$, and take the outcome and constraint spaces to be
}

\draft{

\begin{equation*}
Y \coloneqq \mathbb{R} \times \mathbb{R}^{d} \times \mathbb{R}^{d \times d} \times \mathbb{R}^{d \times d \times d} ,
\end{equation*}

}

\draft{

\begin{equation*}
W \coloneqq \mathbb{R}^{d} \times \mathbb{R}^{d \times d} \times \mathbb{R}^{d \times d \times d} ,
\end{equation*}

}

\draft{
and write $y = \left( y_{0} , \eta , \Theta , \Xi \right) \in Y$. We fix $y_{0} = 1$ and take the
remaining coordinates to lie in the Euclidean unit ball:
}

\draft{

\begin{equation*}
D \coloneqq \left\{ \left( 1 , \eta , \Theta , \Xi \right) : {\left\Vert \eta \right\Vert}_{2}^{2} + {\left\Vert \Theta \right\Vert}_{F}^{2} + {\left\Vert \Xi \right\Vert}_{F}^{2} \le 1 \right\} .
\end{equation*}

}

\draft{
For $x = \left( x_{0} , \xi \right)$, $z \in \mathbb{R}$, and $y \in Y$, define $F \left( x , z , y \right)$ coordinatewise
by
}

\draft{

\begin{equation*}
F_{i} \left( x , z , y \right) \coloneqq z \left( x_{0} \eta_{i} - \frac{1}{2} \xi_{i} y_{0} \right) ,
\end{equation*}

}

\draft{

\begin{equation*}
F_{i j} \left( x , z , y \right) \coloneqq z \left( x_{0} \Theta_{i j} - \xi_{i} \eta_{j} \right) ,
\end{equation*}

}

\draft{

\begin{equation*}
F_{i j \ell} \left( x , z , y \right) \coloneqq z \left( x_{0} \Xi_{i j \ell} - \xi_{i} \Theta_{j \ell} \right) .
\end{equation*}

}

\draft{
Each term contains one coordinate from each of $x$, $z$, and $y$, so $F$ is
trilinear. Define the reward by
}

\draft{

\begin{equation*}
r \left( x , y \right) \coloneqq \frac{1}{m} \sum_{\ell = 1}^{m} \Xi_{i_{\ell} , j_{\ell} , n + \ell} .
\end{equation*}

}

\draft{
It is linear and independent of $x$; its coefficient vector has Euclidean
norm $\frac{1}{\sqrt{m}} \le 1$.
}

\draft{
Fix $x \in X$ and $z \in Z$. Since $z x_{0} \ne 0$ and $y_{0} = 1$ on $D$, we can solve
$F \left( x , z , y \right) = 0$ first for $\eta$, then for $\Theta$, and then for $\Xi$. This gives
}

\draft{

\begin{equation*}
\eta_{i} = \frac{1}{2} q_{i} , \quad \Theta_{i j} = \frac{1}{2} q_{i} q_{j} , \quad \Xi_{i j \ell} = \frac{1}{2} q_{i} q_{j} q_{\ell} .
\end{equation*}

}

\draft{
Writing $t \coloneqq {\left\Vert q \right\Vert}_{2} \le \frac{3}{4}$, the squared norm of $\left( \eta , \Theta , \Xi \right)$ is
}

\draft{

\begin{equation*}
\frac{1}{4} \left( t^{2} + t^{4} + t^{6} \right) \le \frac{4329}{16384} < 1 .
\end{equation*}

}

\draft{
Thus this point lies in $D$, so $K_{z} \left( x \right)$ is a singleton independent of $z$.
The map $F_{x , z} : Y \to W$ is also onto: for any right-hand side, set $y_{0} = 0$
and solve successively for $\eta$, $\Theta$, and $\Xi$.
}

\draft{
Within the hyperplane $y_{0} = 1$, the equations have just one solution, and
that solution lies in $D$. So there are no solutions outside $D$ to check
in the non-tangency condition. We can also prove the distance inequality
for points in the full kernel $L_{z} \left( x \right)$. Since $F_{x , z}$ is onto and the
dimension of $Y$ is one more than that of $W$, this kernel is the line
spanned by $y \left( x \right) = \left( 1 , g \left( q \right) \right)$, where $g \left( q \right) \coloneqq \left( \eta , \Theta , \Xi \right)$ is given by the
equations above. If $p = \tau y \left( x \right) \in L_{z} \left( x \right)$ lies outside $D$, then
}

\draft{

\begin{equation*}
\operatorname{dist} \left( p , D \right) \ge \left| \tau - 1 \right| , \quad \operatorname{dist} \left( p , K_{z} \left( x \right) \right) = \left| \tau - 1 \right| \sqrt{1 + {\left\Vert g \left( q \right) \right\Vert}_{2}^{2}} .
\end{equation*}

}

\draft{
Thus the non-tangency inequality holds for all these points with
}

\draft{

\begin{equation*}
S \ge \frac{1}{\sqrt{1 + \frac{4329}{16384}}} = \frac{128}{\sqrt{20713}} > 0.88 .
\end{equation*}

}

\draft{
Evaluating the reward at the unique point in $K_{z} \left( x \right)$ gives
}

\draft{

\begin{equation*}
v_{z} \left( x \right) = \frac{1}{2 m} p_{G} \left( q \right) .
\end{equation*}

}

\draft{
Since $q$ ranges over the radius-$\frac{3}{4}$ ball, homogeneity and
Lemma~\ref{lem:clique-cubic} imply, for every $z \in Z$,
}

\draft{

\begin{equation*}
V_{z} = \frac{27}{128 m} \operatorname*{max}_{{\left\Vert h \right\Vert}_{2} = 1} p_{G} \left( h \right) = \frac{27}{128 m} \sqrt{\frac{2}{27}} \sqrt{1 - \frac{1}{\omega \left( G \right)}} .
\end{equation*}

}

\draft{
For $j = 2 , \ldots , n$, let
}

\draft{

\begin{equation*}
U_{j} \coloneqq \frac{27}{128 m} \sqrt{\frac{2}{27}} \sqrt{1 - \frac{1}{j}} .
\end{equation*}

}

\draft{
If $\omega \left( G \right) \ge k$, then $V_{z} \ge U_{k}$; if $\omega \left( G \right) \le k - 1$, then
$V_{z} \le U_{k - 1}$. The gap between these bounds satisfies
}

\draft{

\begin{equation*}
U_{k} - U_{k - 1} \ge \frac{27}{128 m} \sqrt{\frac{2}{27}} \frac{1}{2 k \left( k - 1 \right)} = \Omega \left( \frac{1}{m k^{2}} \right) .
\end{equation*}

}

\draft{
Since $\sqrt{\frac{2}{27}} > \frac{1}{4}$ and $m \le n^{2}$, the gap $\Delta_{k} \coloneqq U_{k} - U_{k - 1}$
satisfies
}

\draft{

\begin{equation*}
\Delta_{k} > \frac{1}{40 m k \left( k - 1 \right)} \ge \frac{1}{40 m n^{2}} .
\end{equation*}

}

\draft{
Set $\varepsilon \coloneqq \frac{1}{500 m n^{2}}$. Using polynomially many bits, we can compute a
rational $\tau_{k}$ such that
}

\draft{

\begin{equation*}
\left| \tau_{k} - \frac{1}{2} \left( U_{k} + U_{k - 1} \right) \right| \le \varepsilon .
\end{equation*}

}

\draft{
An estimate of $V_{z}$ with additive error at most $\varepsilon$ lies above
$\tau_{k}$ when $\omega \left( G \right) \ge k$ and below it when $\omega \left( G \right) \le k - 1$, because
$2 \varepsilon < \frac{\Delta_{k}}{2}$. It therefore decides whether $G$ contains a $k$-clique.
The constructed instance has dimension $O \left( d^{3} \right)$, and all its defining
coefficients are rational with polynomial encoding length. This proves the
planning claim.
}

\draft{
Given any rational arm, we can compute the unique point in $K_{z} \left( x \right)$
exactly in polynomial time using the equations above; $x_{0} \ge \frac{2}{3}$ ensures
that we can divide by $x_{0}$. We can therefore simulate nature by returning
this point. Applying Lemma~\ref{lem:planning-from-learning} with the chosen $\varepsilon$
gives an estimate of $V_{z}$ with additive error at most $\varepsilon$, with
probability at least $\frac{2}{3}$. Comparing it with $\tau_{k}$ decides CLIQUE, so
$\text{\normalfont NP} \subseteq \text{\normalfont BPP}$. For a deterministic learner, the estimate is
deterministic and gives
$\text{\normalfont P} = \text{\normalfont NP}$.
}
\end{proof}
\endgroup
}




\draft{\subsection{Proof of Theorem~\ref{thm:iucb-np-hardness}}\label{app:hardness-iucb}}



\draft{
Suppose we are given rational descriptions of $X$, $D$, $Z$, the affine
reward $r$, and the maps defining $K_{z} \left( x \right)$. For each hypothesis $z$, let
}

\draft{

\begin{equation*}
v_{z} \left( x \right) \coloneqq \operatorname*{min}_{y \in K_{z} \left( x \right)} r \left( x , y \right) .
\end{equation*}

}

\draft{
On its first round, IUCB's confidence set is all of $Z$, so its optimistic
value is
}

\draft{

\begin{equation*}
V_{\text{\normalfont IUCB}} \coloneqq \operatorname*{max}_{z \in Z , x \in X} v_{z} \left( x \right) .
\end{equation*}

}

\draft{
We will show that computing this value exactly is NP-hard. The same is
true of finding an arm
}

\draft{

\begin{equation*}
x^{\star} \in \operatorname{argmax}_{x \in X} \operatorname*{max}_{z \in Z} v_{z} \left( x \right) .
\end{equation*}

}







\draft{
\begin{proof}
\draft{
Given a rational order-three tensor
$\mathcal{T} \in \mathbb{Q}^{m \times d \times d_{X}}$, we can view it as a bilinear map
$\mathcal{T} : \mathbb{R}^{d} \times \mathbb{R}^{d_{X}} \to \mathbb{R}^{m}$. Its spectral norm is
}

\draft{

\begin{equation*}
{\left\Vert \mathcal{T} \right\Vert}_{\sigma} \coloneqq \operatorname*{max}_{{\left\Vert u \right\Vert}_{2} \le 1 , {\left\Vert x \right\Vert}_{2} \le 1} {\left\Vert \mathcal{T} \left( u , x \right) \right\Vert}_{2} .
\end{equation*}

}

\draft{
Computing this value exactly is NP-hard \citep{hillar2013most}. We construct an
IUCB instance from $\mathcal{T}$ so that an optimal first arm lets us recover
this norm. Let
}

\draft{

\begin{equation*}
H \coloneqq 1 + \sum_{i , j , k} \left| T_{i j k} \right| , \quad \alpha \coloneqq \frac{1}{2 H} , \quad X \coloneqq \left\{ x : {\left\Vert x \right\Vert}_{2} \le 1 \right\} .
\end{equation*}

}

\draft{
Write $z = \left( z_{0} , u \right)$ and $y = \left( w , s \right)$, and take
}

\draft{

\begin{equation*}
Z \coloneqq \left\{ \left( z_{0} , u \right) : {\left( z_{0} - \frac{5}{3} \right)}^{2} + {\left\Vert u \right\Vert}_{2}^{2} \le 1 \right\} , \quad D \coloneqq \left\{ \left( w , s \right) : {\left\Vert w \right\Vert}_{2}^{2} + s^{2} \le 1 \right\} .
\end{equation*}

}

\draft{
For $z = \left( z_{0} , u \right)$, define the constraint coefficients in
Equation~\ref{eq:setting-compatible-set} and the reward by
}

\draft{

\begin{equation*}
B_{z} x \coloneqq - \alpha \mathcal{T} \left( u , x \right) , \quad C_{z} \left( w , s \right) \coloneqq z_{0} w , \quad d_{z} \coloneqq 0 , \quad r \left( x , \left( w , s \right) \right) \coloneqq \frac{1}{2} \left( 1 + s \right) .
\end{equation*}

}

\draft{
The feasible outcomes at arm $x$ are then
}

\draft{

\begin{equation*}
K_{z} \left( x \right) = \left\{ \left( w , s \right) \in D : w = \frac{\alpha}{z_{0}} \mathcal{T} \left( u , x \right) \right\} .
\end{equation*}

}

\draft{
The definition of $Z$ gives $z_{0} \ge \frac{2}{3}$ and
}

\draft{

\begin{equation*}
\frac{{\left\Vert u \right\Vert}_{2}}{z_{0}} \le \frac{3}{4} ,
\end{equation*}

}

\draft{
because
$1 - {\left( z_{0} - \frac{5}{3} \right)}^{2} - 9 \frac{z_{0}^{2}}{16} = - \frac{{\left( 15 z_{0} - 16 \right)}^{2}}{144} \le 0$.
Equality holds at $z_{0} = \frac{16}{15}$ and ${\left\Vert u \right\Vert}_{2} = \frac{4}{5}$ in every direction.
Also, ${\left\Vert \mathcal{T} \left( u , x \right) \right\Vert}_{2} \le H {\left\Vert u \right\Vert}_{2} {\left\Vert x \right\Vert}_{2}$, so the constraint fixes
$w$ to a vector of norm at most $3 \alpha \frac{H}{4} = \frac{3}{8}$. Thus every $K_{z} \left( x \right)$ is
nonempty, and the uniform non-tangency condition holds with
$S \ge \frac{\sqrt{55}}{8}$. To minimize the reward, nature chooses the smallest
feasible $s$, namely $s = - \sqrt{1 - {\left\Vert w \right\Vert}_{2}^{2}}$. If we write
$h \left( t \right) \coloneqq \frac{1}{2} \left( 1 - \sqrt{1 - t^{2}} \right)$, then
}

\draft{

\begin{equation*}
v_{z} \left( x \right) = h \left( \frac{\alpha}{z_{0}} {\left\Vert \mathcal{T} \left( u , x \right) \right\Vert}_{2} \right) .
\end{equation*}

}

\draft{
For fixed $z$, define the matrix $M_{u}$ by $M_{u} x \coloneqq \mathcal{T} \left( u , x \right)$. A top
right singular vector of $M_{u}$ maximizes $v_{z} \left( x \right)$, so planning takes
polynomial time when $z$ is known. But IUCB must optimize over both $z$
and $x$. Since $h$ is strictly increasing and $\frac{{\left\Vert u \right\Vert}_{2}}{z_{0}}$ can equal
$\frac{3}{4}$ in every direction,
}

\draft{

\begin{equation*}
V_{\text{\normalfont IUCB}} = h \left( 3 \frac{\alpha}{4} {\left\Vert \mathcal{T} \right\Vert}_{\sigma} \right) .
\end{equation*}

}

\draft{
We can therefore recover the tensor norm from the optimistic value:
}

\draft{

\begin{equation*}
{\left\Vert \mathcal{T} \right\Vert}_{\sigma} = \frac{4}{3 \alpha} \sqrt{1 - {\left( 1 - 2 V_{\text{\normalfont IUCB}} \right)}^{2}} .
\end{equation*}

}

\draft{
If we are given only a globally optimal arm $x^{\star}$, define the matrix
$N_{x}$ by $N_{x} u \coloneqq \mathcal{T} \left( u , x^{\star} \right)$. Computing its largest singular value
gives
}

\draft{

\begin{equation*}
\operatorname*{max}_{z \in Z} v_{z} \left( x^{\star} \right) = h \left( 3 \frac{\alpha}{4} \sigma_{\operatorname*{max}} \left( N_{x} \right) \right) = V_{\text{\normalfont IUCB}} .
\end{equation*}

}

\draft{
The same formula then recovers ${\left\Vert \mathcal{T} \right\Vert}_{\sigma}$, so finding a globally
optimal first arm is NP-hard as well.
}
\end{proof}
}



\draft{\subsection{Proof of Theorem~\ref{thm:easy-planning-hard-learning}}\label{app:hardness-easy-planning}}



\draft{
An efficient planner for every hypothesis does not always give an
efficient learner, so the converse of Lemma~\ref{lem:planning-from-learning} fails.
We show this with finite arm and hypothesis sets that have short
descriptions. We also allow the reward to be linear in the arm and in the
outcome separately, so it can contain products of arm and outcome
coordinates.
}







\draft{
\begin{proof}
\draft{
Fix $2 \le k \le n$, and let
}

\draft{

\begin{equation*}
E_{n} \coloneqq \left\{ \left\{ i , j \right\} : 1 \le i < j \le n \right\} , \quad m \coloneqq \binom{n}{2} , \quad q \coloneqq \binom{k}{2} .
\end{equation*}

}

\draft{
For every $k$-element set $T \subseteq \left[ n \right]$, define $x^{T} \in \mathbb{R}^{m}$ by
}

\draft{

\begin{equation*}
x_{e}^{T} \coloneqq \begin{cases}
\frac{1}{q} & \text{\normalfont if } e \subseteq T \\
0 & \text{\normalfont otherwise}
\end{cases} , \quad e \in E_{n} ,
\end{equation*}

}

\draft{
and take
}

\draft{

\begin{equation*}
X \coloneqq \left\{ x^{T} : T \subseteq \left[ n \right] , \left| T \right| = k \right\} .
\end{equation*}

}

\draft{
An arm is represented succinctly by the vertex set $T$. For every
$k$-element set $U \subseteq \left[ n \right]$, define $z^{U} \in \mathbb{R}^{1 + m}$ by
}

\draft{

\begin{equation*}
z_{0}^{U} \coloneqq 1 , \quad z_{e}^{U} \coloneqq \begin{cases}
1 & \text{\normalfont if } e \subseteq U \\
0 & \text{\normalfont otherwise}
\end{cases} ,
\end{equation*}

}

\draft{
and let $Z \coloneqq \left\{ z^{U} : U \subseteq \left[ n \right] , \left| U \right| = k \right\}$.
}

\draft{
Write an outcome as $y = \left( y_{0} , g , s \right) \in \mathbb{R}^{1 + 2 m}$ and set
}

\draft{

\begin{equation*}
D \coloneqq \left\{ \left( 1 , g , s \right) : g \in {\left[ 0 , 1 \right]}^{m} , s \in {\left[ - 1 , 1 \right]}^{m} \right\} .
\end{equation*}

}

\draft{
With constraint space $W = \mathbb{R}^{m}$, let $F_{0} \left( x , z \right) \coloneqq 0$ and define
}

\draft{

\begin{equation*}
{\left( F_{1} \left( y , z \right) \right)}_{e} \coloneqq z_{e} \left( g_{e} - y_{0} - s_{e} \right) + z_{0} s_{e} , \quad e \in E_{n} .
\end{equation*}

}

\draft{
For $z \in Z$, set
}

\draft{

\begin{equation*}
K_{z} \left( x \right) \coloneqq \left\{ y \in D : F_{0} \left( x , z \right) + F_{1} \left( y , z \right) = 0 \right\} , \quad h_{z} \left( x \right) \coloneqq \left\{ Q \in \Delta D : \mathbb{E}_{y \sim Q} \left[ y \right] \in K_{z} \left( x \right) \right\} ,
\end{equation*}

}

\draft{
and let $H \coloneqq \left\{ h_{z} : z \in Z \right\}$.
}

\draft{
The maps $F_{0}$ and $F_{1}$ are bilinear in their displayed arguments. Since
$y_{0}$ is an outcome coordinate, we can write the constraint as $C_{z} y = 0$,
with $B_{z} = 0$, $d_{z} = 0$, and $z \mapsto C_{z}$ linear. For $z = z^{U}$,
}

\draft{

\begin{equation*}
{\left( F_{1} \left( y , z^{U} \right) \right)}_{e} = \begin{cases}
g_{e} - y_{0} & \text{\normalfont if } e \subseteq U \\
s_{e} & \text{\normalfont otherwise} .
\end{cases}
\end{equation*}

}

\draft{
Thus, inside $D$, the constraints force $g_{e} = 1$ for every
$e \subseteq U$ and $s_{e} = 0$ for every $e \nsubseteq U$. The remaining
$g_{e}$ coordinates are free to vary in $\left[ 0 , 1 \right]$. The feasible set is nonempty: take
$g_{e} = 1$ on the edges of $U$, $g_{e} = 0$ elsewhere, and $s = 0$. Moreover,
$y \mapsto F_{1} \left( y , z^{U} \right)$ is onto $W$: given $w \in W$, set $y_{0} = 0$, use
$g_{e} = w_{e}$ and $s_{e} = 0$ on the edges of $U$, and use $g_{e} = 0$ and $s_{e} = w_{e}$
elsewhere.
}

\draft{
Define the reward
}

\draft{

\begin{equation*}
r \left( x , \left( y_{0} , g , s \right) \right) \coloneqq \sum_{e \in E_{n}} x_{e} g_{e} .
\end{equation*}

}

\draft{
The learner is given a succinct description of the entire instance
$\mathcal{M} = \left( X , D , H , r \right)$, including $Z$, $F_{0}$, and $F_{1}$; only the true $U$ is hidden.
The reward lies in $\left[ 0 , 1 \right]$ on $X \times D$. For an arm $x^{T}$, nature
minimizes it by setting every graph coordinate not forced by $U$ to zero, so
}

\draft{

\begin{equation*}
v_{z^{U}} \left( x^{T} \right) = \frac{\binom{\left| U \cap T \right|}{2}}{q} .
\end{equation*}

}

\draft{
Thus $V_{z^{U}} = 1$, uniquely attained by $x^{U}$. Given $z^{U}$, an exact planner
returns $x^{U}$ by setting $x_{e}^{U} = \frac{z_{e}^{U}}{q}$ for every $e \in E_{n}$, which takes
$O \left( m \right) = O \left( n^{2} \right)$ time.
}

\draft{
We now show that a learner with the stated regret bound would solve
CLIQUE. Given a graph $G = \left( \left[ n \right] , E \left( G \right) \right)$ and target clique size $k$
\citep{garey1979computers}, let
}

\draft{

\begin{equation*}
g_{e}^{G} \coloneqq \begin{cases}
1 & \text{\normalfont if } e \in E \left( G \right) \\
0 & \text{\normalfont otherwise}
\end{cases} , \quad y^{G} \coloneqq \left( 1 , g^{G} , 0 \right) \in D .
\end{equation*}

}

\draft{
If $U$ is a $k$-clique of $G$, nature can return $y^{G}$ on every round
under hypothesis $z^{U}$: every edge inside $U$ has $g_{e}^{G} = 1$, and all the
$s_{e}$ coordinates are zero. For every candidate $T$,
}

\draft{

\begin{equation*}
r \left( x^{T} , y^{G} \right) = \frac{\left| E \left( G \left[ T \right] \right) \right|}{q} .
\end{equation*}

}

\draft{
This reward is $1$ exactly when $T$ is a $k$-clique, and is at most
$1 - \frac{1}{q}$ otherwise.
}

\draft{
Assume the policy in the theorem exists, and give it the outcome $y^{G}$
on every round. If $G$ contains a $k$-clique $U$, these outcomes are
compatible with $z^{U}$, whose optimal value is $1$. If the policy plays
$x^{T_{1}} , \ldots , x^{T_{N}}$, the total reward it loses relative to this value is
}

\draft{

\begin{equation*}
{\mathcal{L}}_{N} \coloneqq N - \sum_{t = 1}^{N} r \left( x^{T_{t}} , y^{G} \right) .
\end{equation*}

}

\draft{
This quantity is nonnegative. Let $\mathcal{E}$ be the event that none of
$T_{1} , \ldots , T_{N}$ is a clique. On this event, the learner loses at least
$\frac{1}{q}$ on every round, so ${\mathcal{L}}_{N} \ge \frac{N}{q}$. Markov's inequality and the
assumed regret bound give
}

\draft{

\begin{equation*}
\operatorname{Pr} \left( \mathcal{E} \right) \le q \frac{E \left[ {\mathcal{L}}_{N} \right]}{N} \le q P \left( | \mathcal{M} | \right) N^{- \beta} .
\end{equation*}

}

\draft{
Choose
}

\draft{

\begin{equation*}
N \coloneqq \left\lceil {\left( 4 q P \left( | \mathcal{M} | \right) \right)}^{\frac{1}{\beta}} \right\rceil .
\end{equation*}

}

\draft{
The instance description has length polynomial in $n$, so $N$ and the
time needed to simulate the learner are polynomial in the CLIQUE input
length. Check each $k$-set proposed by the learner and accept as soon as
one is a clique in $G$. If $G$ has a $k$-clique, the preceding bound gives
acceptance probability at least $\frac{3}{4}$. If $G$ has none, none of the
proposed sets can pass this check. In that case, the outcomes may be
incompatible with every hypothesis, but the learner must still terminate
in polynomial time on every syntactically valid history. Thus the
simulation also takes polynomial time when $G$ has no $k$-clique.
}

\draft{
This is an RP algorithm for CLIQUE. Since CLIQUE is NP-complete, it implies
$\text{\normalfont NP} = \text{\normalfont RP}$. If the policy is deterministic, the same verified search gives
$\text{\normalfont P} = \text{\normalfont NP}$.
}
\end{proof}
}

\draft{
This example shows that efficient planning alone does not suffice for
efficient learning. It uses exponentially large finite sets $X$ and $Z$
represented by $k$-subsets, a polytope $D$, and the bilinear reward
$r \left( x , y \right) = \sum_{e} x_{e} g_{e}$. It therefore falls outside our setting of Euclidean
balls and jointly affine rewards.
}
